\section{Finite framed groupoids and the small genera}
\label{sec:B}

We use $E(G)$ for the internal edges, $I(G)$ for the individually
labelled legs, and $F(G)$ for all flags.  Thus the manuscript's
$\widehat A$ is $J(G):=E(G)\amalg I(G)$, and its
$\widehat{\vec A}$ is $F(G)$.  The map
\[
 q:F(G)\longrightarrow J(G)
\]
has two-element fibres over internal edges and one-element fibres over
legs.  All products in the first subsection are on the \emph{pinned}
stratum $M_G=\prod_v M_{g(v),F_v}$; we introduce the action of
$\operatorname{Aut}(G)$ explicitly when we pass to mixed groupoids.
The $S$ in cote 146 retains the tangent directions at the legs as well
as the smoothing directions at the nodes.  This convention is visible
both in the first quotient on scan 100 and in scans 103--110.
The primary manuscript sources are \citet{B_Grothendieck146}
and \citet{B_Grothendieck148}; scan numbers below refer to their
PDF pages, with manuscript folios stated where visually checked.

\subsection{The contracted product, including its finite sign action}

\begin{result}{B.1}{Proved}{The framed torsor and its two quotients}
For a pinned stable graph $G$, put
\[
 P_G=\prod_{j\in J(G),\,M_G} U(L_j),\qquad
 L_i=T_{x_i}C\quad(i\in I(G)),\qquad
 L_e=T_{x_h}C_{v(h)}\otimes T_{x_{\bar h}}C_{v(\bar h)}.
\]
Here $U(L)=L^{\times}/\mathbb R_{>0}$ is the oriented unit-ray circle
bundle; no Hermitian metric is required.  We write
$P_G^E=\prod_{e\in E(G)}U(L_e)$.  Choose a finite set $B_G$ in the
pinned stratum, a nonempty finite set of rays above each member of
$B_G$, and saturate the rays under
$\mu_G=\{\pm1\}^{J(G)}$.  Denote the resulting finite set by $D_G$.
When a graph action is required, take $B_G$ invariant and also
saturate $D_G$ under the finite group $\operatorname{Aut}(G)$.
For vertices of genus zero and valence three or four, the groupoid
\[
 \mathcal S_G:=\Pi_1(P_G,D_G)
\]
is a $\mu_G$-free $\mathbb Z^{J(G)}$-groupoid, and its internal lattice
acts freely on every nonempty Hom-set.  Its geometric quotients are
\[
 \mathcal S_G/(\mathbb Z^{I(G)},\mu_2^{I(G)})
       \simeq\Pi_1(P_G^E,\operatorname{im}D_G),
 \qquad
 \mathcal S_G/(\mathbb Z^{J(G)},\mu_G)
       \simeq\Pi_1(M_G,B_G).
 \tag{B.1}
\]
The first groupoid has the node-framed link group of Part A.  The
second has automorphism group $\prod_{v:\,|F_v|=4}F_2$.  After choices
of trivializations, the automorphism group of $\mathcal S_G$ is
\[
 \mathbb Z^{J(G)}\times\prod_{v:\,|F_v|=4}F_2.
 \tag{B.2}
\]
The direct-product splitting in (B.2) is not specified as
$\operatorname{Aut}(G)$-equivariant.
\end{result}
\begin{proof}
For a trivalent rational component the moduli space is a point.  For
four flags its moduli space is $\mathbb P^1\setminus\{0,1,\infty\}$,
which retracts to a graph of rank two.  Each tangent-ray bundle on
that graph is trivial: construct a section on a maximal tree, and
extend it around each remaining edge using the connectedness of
$S^1$.  Each $L_j$ above is a tensor product of pullbacks of such
bundles.  Consequently their product admits a section, and $P_G$ is
isomorphic, after this choice, to
$M_G\times(S^1)^{J(G)}$.  This proves (B.2) and proves injectivity of
the full lattice without an appeal to the notes' claim of freeness.

The finite sign action is multiplication of nonzero rays.  A sign
fixing a ray in every coordinate has all coordinates $+1$, so its
action on $D_G$ is free.  For paths with prescribed endpoints, the
projection to the base has fibres that are torsors under the winding
lattice.  Dividing by the leg-coordinate winding lattice and
identifying the corresponding sign-orbits of endpoints therefore
gives exactly the first groupoid in (B.1); dividing by all coordinates
gives the second.  In particular, the quotient notation includes
both the identification of endpoint lifts and the killing of their
vertical loops.  It does not mean the action groupoid of the finite
sign group.

For completeness, the contraction from flags to edges is the map of
torus bundles
\[
 \prod_{h\in F(G)}U(T_{x_h}C_{v(h)})\longrightarrow P_G,
 \qquad (u_h)_h\longmapsto
 \bigl((u_hu_{\bar h})_{e=\{h,\bar h\}},(u_i)_{i\in I}\bigr).
 \tag{B.3}
\]
Its map on integer lattices sends both $e_h,e_{\bar h}$ to $e_e$,
and its map on signs multiplies the two signs.  Its kernel is the
connected torus with one coordinate per edge, embedded by
$u\mapsto(u,u^{-1})$.  Killing its full fundamental groupoid,
including the paths between its different sign lifts, is the
contracted product in scan 100.  To check the assertion on each
fibre, lift the circles to $\mathbb R^{F(G)}$ and
$\mathbb R^{J(G)}$: the same coordinate-sum map is surjective, and
its kernel is generated by $e_h-e_{\bar h}$.  Two lifted paths have
the same image precisely when their difference is a kernel path.
This proves the universal quotient statement on arrows and
composition.  Applying the construction over the product base
proves (B.3) on fundamental groupoids.  Graph automorphisms permute
flags and edges and commute with (B.3), so these constructions carry
the required graph action before taking mixed quotients.
\end{proof}

\begin{result}{B.2}{Proved}{The presentation actually written in scans 115--117}
For a trivalent rational graph set $r=|J(G)|$.  The finite model of
cote 146, scans 114--117 (folios 113--116), has the $2^r$ objects of
the sign torsor $D_G=\prod_{j\in J(G)}\varepsilon_j$.  It has
$r2^r$ positive half-turn generators
\[
 u_{j,x}:x\longrightarrow\sigma_jx
 \quad(j\in J(G),\ x\in D_G).
\]
Its relations are the source-compatible commuting squares for two
different coordinates.  The commutation of the full-turn loops
$v_{j,x}=u_{j,\sigma_jx}u_{j,x}$ written on scan 117 follows from
these squares.  Equivalently,
\[
 \operatorname{Hom}(x,y)=
 \{m\in\mathbb Z^r:m\bmod2=y-x\},
 \tag{B.4}
\]
where addition gives composition and $m$ records displacement in
units of a half turn.  The loop lattice is $2\mathbb Z^r$, identified
with $\mathbb Z^r$ by $2e_j\mapsto e_j$.
\end{result}
\begin{proof}
Choose a member of the sign torsor only to write its members as
binary vectors.  The universal covering cubulation of the torus has
vertex set $\mathbb Z^r$ and oriented coordinate edges; quotient by
$2\mathbb Z^r$ gives the displayed vertices and generators.  A
sequence of coordinate edges has a displacement $m$.  Commuting
squares interchange adjacent steps in different coordinates, and
backtracking cancels inverse steps.  Sorting coordinate steps and
then cancelling proves that two paths with the same endpoints have
the same class exactly when their displacement vectors agree.
This proves (B.4), with no unproved relation among winding numbers.
The two paths consisting of two steps in coordinate $i$ and two in
coordinate $j$ are interchanged by four commuting squares, proving
the stated redundancy of scan 117's second relation.

A half-turn is therefore an arrow between two distinct chosen
framings.  Its square is a full twist; imposing order two on this
arrow would kill a nonzero member of the meridian lattice.  The
formula also proves the cocycle on scan 116: the discrepancy between
two chosen binary half-turns and their chosen product is one full
turn in precisely the coordinates where both signs are negative.
\end{proof}

\begin{result}{B.3}{Refuted}{The proposed equivalence before forgetting leg framings}
The proposed equivalence $S_0\mathcal T_G\simeq S_A\Pi M_G$ without
quotienting the leg framings is false, already for the graph with
one trivalent rational vertex and three labelled legs.  The exact
comparison is the first functor in (B.1), whose kernel at every
object is $\mathbb Z^{I(G)}$.
\end{result}
\begin{proof}
For this graph $M_G$ is a point and $E(G)=\varnothing$.  Formula
(B.4) gives automorphism group $\mathbb Z^3$ for the scan 115 model;
the early two-object model of scans 3--8 has the same automorphism
group.  The groupoid $S_A\Pi M_G$ has trivial automorphism group.
The forgetful functor sends $(m_1,m_2,m_3)$ to zero and is not
faithful.  No equivalence between the two groupoids exists because
an equivalence induces a bijection on each endomorphism set.  The
quotient by this displayed kernel is the terminal groupoid, which
is the required node-only model.  For general $G$ the same argument
in each fibre gives the comparison (B.1).
\end{proof}

\begin{result}{B.4}{Proved}{The genus-zero tree index and the chart counts}
The category of reduced stable trees with a fixed individually
labelled leaf set, and specified edge contractions fixing those
labels, is equivalent to a finite poset.  On four and five leaves,
the numbers of trees by number of internal edges are respectively
$(1,3)$ and $(1,10,15)$.  For the disjoint union of their trivalent
charts, the scan 116 presentations have the following counts;
these counts precede any global identification of chart objects.
\[
\begin{array}{c|r|r|r|r|r}
 n&\text{charts}&r&\text{framed objects}&\text{half-turn generators}
       &\text{node-only objects / generators}\\\hline
 4&3&5&96&480&6\,/\,6\\
 5&15&7&1920&13440&60\,/\,120
\end{array}
\tag{B.5}
\]
\end{result}
\begin{proof}
Every internal edge separates the labelled leaves into two nonempty
sets, each of cardinality at least two.  Distinct edges give distinct
bipartitions.  The collection of these bipartitions reconstructs the
tree: the vertices are the compatible regions determined by the
edge cuts, and incidence is given by changing one region across one
cut.  Thus an automorphism fixing every leaf fixes every edge and
then every vertex.  A contraction removes precisely its contracted
edge bipartitions.  Its source and target therefore determine the
contracted set uniquely.  This proves both absence of nontrivial
automorphisms and uniqueness of a contraction, the two ingredients
needed for a poset.  A stable tree on $n$ leaves has at most $n-3$
internal edges, so there are finitely many collections of the finite
set of leaf bipartitions.

For four leaves there are three bipartitions into two pairs.  For
five leaves there are ten bipartitions of type $2+3$; two such cuts
are compatible exactly when their two-element parts are disjoint.
There are $\binom52\binom32/2=15$ such pairs.  A trivalent tree has
$n-3$ internal edges, hence $r=2n-3$ framed coordinates.  Apply B.2
and sum over the three or fifteen charts.  The program
\texttt{code/b\_finite\_checks.py} independently enumerates stable
multigraphs, fixes all labels, checks stability, and obtains all
entries of (B.5).  This also specifies the precise finite counting
convention; adding further marked basepoints changes a presentation's
object and generator counts.
\end{proof}

\begin{result}{B.5}{Refuted}{Freeness does not follow from an unspecified finite colimit}
The implication that a finite colimit of groupoids with a free
$\mu_2$-action on objects has a free action on objects is false.
Consequently that implication alone cannot justify the sentence
following the first formula of cote 146, scan 100.
\end{result}
\begin{proof}
Let $C$ be the indiscrete groupoid on the two objects $0,1$, and let
$\mu_2$ interchange them.  The two equivariant functors
$\operatorname{id},\sigma:C\rightrightarrows C$ have coequalizer
with one object.  Its inherited $\mu_2$-action fixes that object.
More explicitly, if $u:0\to1$ is the unique arrow, the coequalizer
identifies $u$ with $u^{-1}$, and has automorphism group $C_2$.
The action on objects is therefore no longer free.  The program
checks the two-object coequalizer.  In B.1 freeness was proved from
the action on actual nonzero rays; it was not deduced from the
existence of a colimit.  The present counterexample addresses the
inference about arbitrary finite colimits, without substituting
an index category for the unfinished index on scan 100.
\end{proof}

\begin{result}{B.6}{Proved}{An explicit obstruction for the boundary-incidence interpretation}
Let $L$ be the boundary of a sufficiently small regular
neighbourhood of the boundary divisor in $\overline M_{0,5}$.  The
colimit of the fundamental groupoids of its circle-bundle pieces
and all their torus intersections presents $\Pi_1(L)$.  The inclusion
$L\to M_{0,5}$ does not induce an isomorphism of fundamental groups.
Thus this specified boundary-incidence colimit cannot be the
missing global presentation in scan 100.
\end{result}
\begin{proof}
There are ten rational boundary components, each meeting three
others, and fifteen intersections.  Their incidence graph has the
ten two-element subsets of a five-element set as vertices and
joins disjoint subsets.  It is the Petersen graph.  Removing small
discs around the intersections leaves ten pairs of pants.  Above
each is a circle bundle, necessarily trivial since a pair of pants
has the homotopy type of a graph.  Thus each vertex group is
$\mathbb Z\times F_2$, each intersection group is $\mathbb Z^2$,
and both edge inclusions are injective: one generator is the circle
fibre, and the other is an infinite-order peripheral loop in the
pair of pants.  Van Kampen with basepoints on all intersections
proves the asserted colimit.

The universal cover obtained by developing this graph of spaces is
a tree of contractible universal covers of these pieces, glued
along contractible universal covers of the tori.  Its finite
subtrees are contractible by successive gluing along contractible
subspaces, and every compact sphere lies in a finite subtree.
Thus $L$ is aspherical.  It is a connected closed oriented
three-manifold, so
\[
 H_3(\pi_1L;\mathbb Z)=H_3(L;\mathbb Z)=\mathbb Z.
\]
On the other hand, the map
$M_{0,5}\to M_{0,4}$ forgetting the fifth point has fibre a sphere
minus four points.  Both base and fibre are graphs up to homotopy,
of ranks two and three.  A section can be given in affine
coordinates $(x,y)$ by $y=2+|x|$.
The resulting extension is $F_3\rtimes F_2$.  Its classifying space
is a graph of mapping cylinders of homotopy equivalences of a
rank-three graph, one cylinder for each free base generator.
Its universal cover is a tree of contractible graph covers and
mapping cylinders, so it is a $K(\pi,1)$ of dimension two.  Hence its third homology is zero.  An isomorphism induced by
$L\to M_{0,5}$ would identify these third homology groups, a
contradiction.  This proof specifies the actual comparison map and
the invariant that prevents it from being an equivalence; it makes
no assertion that the manuscript had defined this index category.
\end{proof}

\begin{result}{B.7}{Known}{The precise relation to the Lego groupoid}
Let $\mathcal L(\Sigma)$ be the marking two-complex of
Bakalov--Kirillov with their moves and relations, and let $\Gamma$
be the mapping class group with the same orientation and boundary
conventions.  The groupoid with objects the markings and arrows
\[
 M\longrightarrow M':\quad (g,[p]),\qquad
 g\in\Gamma,\quad p:gM\leadsto M',
\]
is equivalent to $B\Gamma$ by $(g,[p])\mapsto g$.  With fixed labels
and capped boundary circles it gives the open groupoid
$B\Gamma_{0,n}$ in Part A.  The capping functor kills exactly the
central lattice of boundary twists.
\end{result}
\begin{proof}
Theorem 4.9 of \citet{B_BakalovKirillov} states connectedness and
simple connectedness of $\mathcal L(\Sigma)$ in genus zero; its
proof in \S6 reduces to the sphere mapping-class presentation in
Proposition 6.6 and the commuting lifts of cutting moves.  Therefore
for any $g,M,M'$ there is exactly one relative-endpoint path class
$p$.  The displayed map is bijective on every Hom-set and
essentially surjective.  Composition is
\[
 (h,[q])\circ(g,[p])=(hg,[q*(hp)]),
\]
where the path first traverses $hp$.  This proves the equivalence
with its exact composition map.  Capping sends a homeomorphism
fixed on every boundary to a homeomorphism fixing the new labelled
punctures.  A puncture-fixing isotopy can be adjusted near each
puncture to preserve a small circle; its residual winding is an
integer twist in that collar.  Conversely those collar twists
become isotopic to the identity after capping.  This proves both
surjectivity and the stated kernel.  It is the same quotient as the
leg quotient in B.1.  These maps give a comparison of the groupoids.
They do not assign a word to each arrow in the different, unfinished
finite diagram of cote 148; those words require a chosen set of
geometric paths for that diagram.
\end{proof}

\begin{result}{B.8}{Open}{The remaining sphere-only and original-index comparison}
B.27 proves a strict finite geometric colimit with the five
pieces $(0,3),(0,4),(0,5),(1,1),(1,2)$, including all
transition functors, their coherences, a comparison inverse,
and a proof of the free external sign action.  In genus zero
its index is the finite reduced labelled-tree contraction poset.
B.8a also gives an abstract finite colimit using the central
extensions of the three- and four-point sphere groups.

The remaining comparison is to identify a sphere-only geometric
diagram of the original scan100, including its elliptic cut-changing
moves and their framed central relations, with those proved
constructions, generator by generator.  The bare index phrase
in scan100 does not specify that additional dictionary.
The early four-point model itself is completely reconstructed
in B.16--B.17.  The full five-point realization B.25 is explicit,
with the path-convention comparison isolated in B.15.
\end{result}
The maps in B.1, the non-preservation example B.5, the chosen
boundary-incidence map B.6, and the marking equivalence B.7
retain their stated scopes.  B.27 supplies an actual finite
geometric recipe; the open row here concerns the stronger
sphere-only identification, not existence of that recipe.

\subsection{A finite colimit whose maps are fully specified}

\begin{result}{B.8a}{Proved}{Finite presentation yields a finite colimit of elementary framed groupoids}
For each pinned stable graph, its framed fundamental groupoid admits
a finite-colimit presentation in the category of groupoids with its
specified central lattice and finite sign action.  One may take the
elementary automorphism groups, after changing the central lattice,
to be $\mathbb Z^r$ and $\mathbb Z^r\times F_2$, the groups of the
three- and four-point rational pieces.  The arrows in this colimit
are the word maps specified below.  This proves existence of a
finite algebraic presentation of the asserted kind.  It does not
identify these word maps with the unfinished tree-indexed diagram
of scan 100.
\end{result}
\begin{proof}
Put $\pi=\mathbb Z^r$ for the framed coordinates and let $H$ be the
framed group.  The quotient $\Gamma=H/\pi$ is finitely presented.
For genus zero this follows from the braid presentation and capping;
for positive genus it follows from
\citet[Theorems 3.1 and 3.2]{B_LabruereParis}, followed by the finite
Schreier rewriting for the subgroup fixing all labels.  A product
of vertex groups is finitely presented; gluing two boundary
coordinates imposes one central relation per edge.  Thus the same
assertion holds for every pinned stable graph.

Write $\Gamma=\langle x_1,\ldots,x_m\mid R_1,\ldots,R_s\rangle$
and choose lifts $\widetilde x_i\in H$.  Each lifted relator equals
a uniquely determined element $c_j\in\pi$.  Then
\[
 H=\left(\pi\times F(x_1,\ldots,x_m)\right)
       /\langle\!\langle R_jc_j^{-1}:1\leq j\leq s\rangle\!\rangle.
 \tag{B.6}
\]
To verify injectivity of this presentation into $H$, take a word
mapping to the identity.  Its image in $\Gamma$ is a product of
conjugates of the $R_j^{\pm1}$.  Replace these relators by their
specified central values using (B.6).  What remains lies in $\pi$
and maps to zero in its injective copy in $H$; hence it is zero.
Surjectivity follows by lifting a word in the $x_i$ and adjusting
its residual element of $\pi$.

Let $E=\tfrac12\mathbb Z^r$, with the inclusion
$\pi\subset E$, and form $H_0=H\amalg_{\pi}E$ with $E$ central.
There is a homomorphism $\chi:H_0\to E/\pi=\mu_2^r$ that kills the
chosen $\widetilde x_i$ and restricts to the quotient on $E$.
The translation groupoid of $H_0$ acting on $\mu_2^r$ by $\chi$ is
the framed groupoid on one full sign orbit; its vertex group is
$H$.  Formula (B.6), now with $E$ in place of $\pi$, gives its
finite-colimit construction explicitly: for each relator use the
parallel maps
\[
 E\times F(t)\rightrightarrows E\times F(x_1,\ldots,x_m),
 \qquad t\mapsto R_jc_j^{-1},\quad t\mapsto1,
 \tag{B.7}
\]
with both maps the identity on $E$.  Their coequalizer imposes
exactly (B.6).  The group $E\times F_m$ is a finite amalgamated
coproduct over $E$ of its one-generator free factors.  The
one-generator factor itself is the coequalizer of the identity
and the map $t\mapsto t,u\mapsto1$ on $E\times F(t,u)$.
Consequently every group appearing as a piece can be replaced by
$E$ or $E\times F_2$ in a finite diagram.  All maps preserve
$\chi$ and induce functors that are the identity on the sign-object
set.  Colimits on arrows therefore give the translation groupoid
of $H_0$, with its free sign action.  The proof verifies the
particular diagram's freeness and avoids the false general
inference of B.5.  Additional finitely many basepoints are added
by chosen connecting isomorphisms, which also constitute a finite
diagram.
\end{proof}

\subsection{Explicit presentations in types $(1,1)$ and $(1,2)$}

In the next statements $a,b,c$ are positive Dehn twists on a
three-chain: consecutive curves intersect once and the first and
third are disjoint.  With one marked point we use only $a,b$.
The choices consist of an oriented reference surface, these curves,
and paths from the chosen framed objects to this reference.  The
following homomorphisms are genuine homomorphisms for those choices.
Changing a reference path conjugates the corresponding
homomorphism by its endpoint path; we do not call the resulting
based isomorphisms canonical.

\begin{result}{B.9}{Proved}{The two open groupoids and their leg-framed extensions}
For complex curves with individually labelled markings,
\[
 \Gamma_{1,1}=\langle a,b\mid aba=bab,\ (ab)^6=1\rangle
                 \cong\operatorname{SL}_2(\mathbb Z),
 \tag{B.8}
\]
\[
 \Gamma_{1,2}=\langle a,b,c\mid aba=bab,\ bcb=cbc,
               \ ac=ca,\ (abc)^4=1\rangle.
 \tag{B.9}
\]
Thus $\Pi_1M_{1,n}$ is equivalent to the one-object groupoid on the
displayed group.  The groups of the unit tangent-ray bundles at
all marked points are respectively
\[
 H_{1,1}=\langle a,b\mid aba=bab\rangle,
       \qquad z_1=(ab)^6,
 \tag{B.10}
\]
and
\[
 H_{1,2}=\left\langle a,b,c,z_1,z_2\ \middle|\
 \begin{array}{l}
 aba=bab,\ bcb=cbc,\ ac=ca,\\
 z_1,z_2\text{ central},\ (abc)^4=z_1z_2
 \end{array}\right\rangle
 \cong B_4\times\mathbb Z.
 \tag{B.11}
\]
Here $z_i$ is the positive full turn at the $i$th marked point.
The capping maps are the quotients $z_i\mapsto1$.
\end{result}
\begin{proof}
For two labelled points $p,q$ on a smooth genus-one curve, the
complete linear system $|p+q|$ gives a double cover of a projective
line, with $p,q$ over infinity and four distinct finite branch
points.  An affine coordinate $x$ on the line is unique up to
$x\mapsto\alpha x+\beta$.  Ordering $p,q$ fixes the sign of $y$
at infinity in a monic equation
\[
 y^2=\prod_{i=1}^4(x-r_i).
\]
Every isomorphism of the pointed curves induces the affine map of
the quotient and has exactly one lift preserving the ordering of
$p,q$.  Conversely that lift gives the pointed-curve isomorphism.
These constructions apply in families, giving the analytic-stack
equivalence
\[
 M_{1,2}\simeq[\operatorname{UConf}_4(\mathbb C)/
                        \operatorname{Aff}(\mathbb C)].
\]
The unordered configuration space has braid group $B_4$, with its
usual adjacent-crossing generators and relations.  The translation
subgroup of the affine group is contractible; the positive loop in
its dilation subgroup rotates all four roots once, giving the
central braid $(abc)^4$.  The homotopy exact sequence for this
quotient stack gives exactly (B.9).  The full rotation injects into
$B_4$, since its image under exponent sum is $12$.

At either point over infinity, $u=1/x$ is a local coordinate.
Under dilation it is multiplied by $\alpha^{-1}$.  Keeping two
nonzero tangent vectors replaces the preceding quotient by the
quotient of $\operatorname{UConf}_4(\mathbb C)\times(\mathbb C^*)^2$
with these two weights.  The affine loop therefore imposes exactly
$(abc)^4z_1^{-1}z_2^{-1}=1$.  This proves (B.11), including the
identification of the two central generators and the capping map.
Eliminating $z_2$ exhibits $B_4\times\mathbb Z$.

For one labelled point the double cover from $|2p|$ has infinity
ramified and three finite branch points.  In a monic cubic equation
the allowed scaling is $x\mapsto\lambda^2x+\beta$,
$y\mapsto\lambda^3y$.  A positive loop of $\lambda$ rotates the
three roots twice, giving $(ab)^6$.  The local parameter $x/y$ at
$p$ has weight $\lambda^{-1}$.  The same calculation proves
(B.8) and (B.10).  Mapping $a,b$ to
\[
 \begin{pmatrix}1&1\\0&1\end{pmatrix},\qquad
 \begin{pmatrix}1&0\\-1&1\end{pmatrix}
\]
identifies (B.8) with $\operatorname{SL}_2(\mathbb Z)$: the braid
presentation becomes the standard amalgam
$C_4*_{C_2}C_6$ on $aba$ and $ab$.
The same formulas are obtained by specializing
\citet[Theorem 3.1 and Proposition 2.12]{B_LabruereParis} to
one and two boundary components.  In the two-boundary case their
Artin diagram is the chain $a,b,c$ together with one central
boundary generator; their chain relation identifies the other
boundary generator with $(abc)^4z_1^{-1}$.
\end{proof}

\begin{result}{B.10}{Proved}{All boundary graph groups in these two types}
For $(1,1)$ there is one boundary graph: a rational vertex with one
loop and one labelled leg.  Its pinned stratum is a point and its
mixed stratum is $BC_2$.  Put $h=(ab)^3$.  Its fully framed mixed
link group is
\[
 \langle a,h\mid[a,h]=1\rangle,\qquad z_1=h^2.
 \tag{B.12}
\]
Its node-only mixed link group is
$\langle a,h\mid[a,h]=1,h^2=1\rangle$.

For $(1,2)$ name the four boundary graph types as follows:
$N$ is a rational vertex with a loop and both legs;
$T$ is an elliptic vertex joined to a rational tail carrying both
legs; $C$ is a two-edge rational cycle with one labelled leg on
each vertex; and $D$ is a rational loop vertex joined to a rational
tail carrying both legs.  Write
\[
 s=aba,\quad\Delta=abcaba,\quad
 r=s^{-1}\Delta s,\quad q=(abc)^2,\quad h=(ab)^3,
 \quad d=h^2.
 \tag{B.13}
\]
The pinned stratum groups are $F_2$, $\operatorname{SL}_2(\mathbb Z)$,
$1$, $1$ for $N,T,C,D$, respectively.  The mixed stratum groups are
$C_2*\mathbb Z$, $\operatorname{SL}_2(\mathbb Z)$, $C_2$, $C_2$.
The following are presentations of their fully framed mixed link
groups; in every row $z_1,z_2$ are central:
\[
\begin{array}{c|l|l}
 \text{graph}&\text{generators and additional relations}
                   &\text{node twists}\\\hline
 N&a,c,r,z_1,z_2;\ [a,c]=[a,r]=1,\ r^2=z_1z_2&a\\
 T&a,b,z_1,z_2;\ aba=bab&d=(ab)^6\\
 C&a,c,q,z_1,z_2;\ [a,c]=1,\ qaq^{-1}=c,\ q^2=z_1z_2&a,c\\
 D&a,h,z_1,z_2;\ [a,h]=1&a,d=h^2
\end{array}
\tag{B.14}
\]
The pinned full-framing subgroups are the kernels of the graph
permutation: in rows $N,C,D$ the respective quotient maps send
$r,q,h$ to the nonidentity member of $C_2$, and all lattice
coordinates to the identity.  The graph group in row $T$ is
trivial.  The node-only link groups are obtained by $z_1=z_2=1$.
The mixed stratum groups are obtained by also killing the listed
node twists.
\end{result}
\begin{proof}
The graph enumeration is finite.  For one marking it gives the
open elliptic graph and the one-loop graph.  For two markings it
gives the open graph and exactly $N,T,C,D$.  The program enumerates
all connected multigraphs with at most $3g-3+n$ edges, assigns vertex
genera with total genus $g$, assigns the individual legs, and checks
$2g(v)-2+\operatorname{val}(v)>0$.  It obtains graph automorphism
orders $2,1,2,2$ for $N,T,C,D$, as asserted.

Cutting the node of $N$ leaves a four-point sphere.  Pin the two
external markings at $0,\infty$ and the two node branches at $1,z$.
Exchanging the branches acts by $z\mapsto z^{-1}$.  Its quotient has
coordinate $w=z+z^{-1}$, two missing points $2,\infty$, and a single
orbifold point of order two at $-2$.  Its orbifold fundamental group
is therefore $\langle r,c\mid r^2=1\rangle=C_2*\mathbb Z$.
The subgroup fixing the branches is free on $c,rcr^{-1}$.  The
node circle supplies a commuting generator $a$.  At the order-two
configuration, the involution of the four-point sphere is
$x\mapsto-x$: its derivatives at the two external markings are
$-1$, and the product of its derivatives at the two branches is
$+1$.  A framed lift consequently has square $z_1z_2$, with no
node-winding factor.  This proves row $N$ by the presentation of a
central extension of $C_2*\mathbb Z$; there is no additional
relator on its free generator $c$.

Cutting the separating node of $T$ leaves a one-point elliptic
component and a rigid rational component.  Keeping the node
framing converts the elliptic factor to $B_3$ by (B.10); the two
external leg framings supply $z_1,z_2$.  Its boundary twist is
$(ab)^6$.  This proves row $T$.

For $C$ both rational components are rigid.  The node and leg
lattices have rank four.  The graph involution interchanges the two
nodes, and its lift $q$ interchanges $a,c$; its square is $z_1z_2$
by (B.11).  The index-two subgroup generated by
$a,c,z_1,z_2$ is free abelian: its generators are independent twists
about the four boundary or cutting circles.  Cutting into annuli
and pairs of pants proves that a product of these twists is
trivial only when every exponent is zero.  Thus the extension by
$C_2$ has precisely row $C$'s presentation.

For $D$, the involution flips the branches of the rational loop.
After smoothing the loop it is the elliptic half-turn $h$ in the
one-holed torus, whose square is the separating-node twist $d$.
It commutes with $a$ and both external twists.  The subgroup on
$a,h^2,z_1,z_2$ is again the independent rank-four lattice, so the
displayed extension has no further relation.  The same argument
with one external marking gives (B.12).  All passages to pinned
and node-only groups are the specified kernel and quotient maps;
the finite graph actions have not been discarded.
\end{proof}

\begin{result}{B.11}{Proved}{Generization on every boundary flag in the two small genera}
The generization maps from the groups (B.12)--(B.14) to the open
framed groups are the inclusions specified by the words (B.13).
In $(1,2)$ the two nontrivial two-step flags, in addition to the
cycle flags, have maps
\[
 D\longrightarrow T:\quad a\mapsto a,\ h\mapsto(ab)^3,
 \qquad
 D\longrightarrow N:\quad a\mapsto a,\ h\mapsto ac^{-1}r,
 \tag{B.15}
\]
with $z_i\mapsto z_i$.  For $C\to N$ retaining the node $a$, the
arrow-stabilizer group is the subgroup
$\langle a,c,z_1,z_2\rangle$ of row $C$, and its map to row $N$ is
the identity on these named generators.  The other contraction is
its conjugate by $q$.  These maps compose to the maximal
generization maps.  After the lattice quotients in B.10 they give
the corresponding pinned and mixed groupoid spans of Part A.
\end{result}
\begin{proof}
A positive rotation of a plumbing parameter is the positive Dehn
twist about its vanishing circle.  Therefore each lattice
coordinate has the stated image.  The remaining generators are
homeomorphisms of the cut components extended over the annuli;
this gives the indicated homomorphisms into the stabilizers of the
remaining curves.  The braid relations give
\[
 \Delta^2=(abc)^4,\quad rar^{-1}=a,\quad
 q a q^{-1}=c,\quad q c q^{-1}=a.
\]
They also give the exact word identity needed in (B.15):
\[
 \begin{split}
 ac^{-1}r
 &=ac^{-1}a^{-1}b^{-1}a^{-1}(abcaba)(aba)\\
 &=(aba)^2=(ab)^3=h.
 \end{split}
\]
Thus both routes from $D$ to the open group agree as genuine based
homomorphisms.  The same is true of the cycle maps on their
arrow-stabilizers.  The element $q$ exchanges the two cycle edges,
so it does not preserve the choice of the one edge being smoothed;
this is exactly why it is absent from that arrow-stabilizer.  There
is no map from all of row $C$ to that chosen branch of row $N$
asserted here.  Including its second contraction gives the other
branch and the required graph equivariance.

For verification the program \texttt{b\_braid\_checks.py} evaluates
both sides of all the displayed braid identities in the faithful
Artin action on a free group.  Its twelve tests all pass.  The
proof of faithfulness is not needed to infer the homomorphisms,
which follow from the geometric construction; the computation is
an independent exact check of the words used in that construction.
The complete lists of boundary graphs and contractions show that
no further flags occur in these two types.
\end{proof}

\subsection{The finite five-point diagram, checked against its original image}

\begin{result}{B.12}{Proved}{The nine vertex orbits and the corrected sixteen arrow orbits}
Let $K$ be a fixed five-element set and let $\mathfrak S_K$ act by
relabeling.  The vertex sets of cote 148, scan 31, are the following
nine transitive $\mathfrak S_K$-sets.  The symbols in the first column
are the manuscript's symbols; subscripts are stabilizer orders.
\[
\begin{array}{c|c|r}
\text{type}&\text{stabilizer in }S_5&\text{number}\\ \hline
\mathrm I_{10}&D_5&12\\
\mathrm I_6&S_3\text{ acting on }3+2\text{ by its sign}&20\\
\mathrm I_4&C_4&30\\
\mathrm {II}_2&C_2&60\\
\mathrm {II}'_2&C_2&60\\
\mathrm {II}_1&1&120\\
\mathrm {III}_2&C_2&60\\
\mathrm {III}_1&1&120\\
\mathrm {III}'_2&C_2&60
\end{array}
\tag{B.16}
\]
In every $C_2$ row the nonidentity element is a double
transposition.  Thus the number of vertices is $542$, in nine
orbits.  The fourteen named arrow families have the following
cardinalities and orbit decompositions:
\[
\begin{array}{c|c|r|r}
\text{family}&\text{source and target}&\text{arrows}&S_5\text{-orbits}\\ \hline
 f_{B,s}&\mathrm I_6\to\mathrm I_4&60&1\\
 f_{\vec Q,a}&\mathrm {II}'_2\to\mathrm I_4&120&1\\
 \tau_{\pi,s}&\mathrm {II}_2\to\mathrm {II}'_2&60&1\\
 \tau_r^0&\mathrm {II}_1\to\mathrm {II}_1&120&1\\
 \gamma_{\pi,s}&\mathrm {II}_2\to\mathrm I_{10}&60&1\\
 \gamma''_{\pi,s}&\mathrm {II}'_2\to\mathrm I_6&60&1\\
 \gamma_r^0&\mathrm {II}_1\to\mathrm I_6&120&1\\
 \tau_{\pi,\rho}&\mathrm {III}_2\to\mathrm {III}_1&120&1\\
 \gamma_{\pi,\rho}&\mathrm {III}_2\to\mathrm {II}_2&120&1\\
 \tau_r^1&\mathrm {III}_1\to\mathrm {III}'_2&120&1\\
 \gamma_r^1&\mathrm {III}_1\to\mathrm {II}'_2&120&1\\
 \gamma_r^{1'}&\mathrm {III}_1\to\mathrm {II}_1&360&3\\
 \gamma_r^2&\mathrm {III}'_2\to\mathrm {II}_2&120&1\\
 \gamma_r^{2'}&\mathrm {III}'_2\to\mathrm I_4&60&1
\end{array}
\tag{B.17}
\]
Consequently $1620$ arrows means sixteen orbits, with orbit-size
histogram $5\cdot60+11\cdot120$.  The reported version ``fourteen
orbits'' is false; fourteen counts the named families in (B.17).
The original scan already strikes out the word \emph{orbites} next
to $14$, so this correction is not attributed to an unqualified
retained assertion of the author.
\end{result}
\begin{proof}
An unoriented pentagonal order on $K$ has dihedral stabilizer of
order ten, giving $120/10=12$ objects of type $\mathrm I_{10}$.
A bitetrahedral structure is a partition $K=J\amalg\varepsilon$,
$|J|=3$, together with a bijection from $\varepsilon$ to the two
orientations of $J$.  There are $\binom53\cdot2=20$.  Its stabilizer
permutes $J$ arbitrarily and acts on $\varepsilon$ by the sign of
that permutation.  An oriented square on $K\setminus\{P\}$ has
cyclic stabilizer of order four, and there are $5\cdot6=30$ such
objects.  A pentagon with a distinguished vertex has stabilizer
the reflection fixing that vertex, a double transposition; there
are $12\cdot5=60$.  This describes the two $\mathrm {II}$ types
with subscript two and $\mathrm {III}_2$.  A pinned structure of
$\mathrm {II}_1$ has data
\[
 K=J\amalg\varepsilon,\quad |J|=3,\quad
 \omega\in\operatorname{Or}(J),\quad r\in J\times\varepsilon.
\]
There are $10\cdot2\cdot6=120$, and fixing all these data fixes
each label.  The fully pinned $\mathrm {III}_1$ structures have
the same trivial stabilizer.  The square with a specified
codiagonal frame at $\mathrm {III}'_2$ has a remaining half-turn,
and hence $120/2=60$ objects.  These constructions prove (B.16)
as identifications of $S_5$-sets, not just as a sum of integers.

For the arrows, the first family has parameters $(B,s)$ with
$s\in J(B)$; it has $20\cdot3=60$ members and double-transposition
stabilizer.  The second has an oriented square and one of its four
specified flags, hence $30\cdot4=120$ members and trivial
stabilizer.  The families parameterized by $(\pi,s)$ have the
pentagon-with-vertex stabilizer.  The families parameterized by a
complete frame $r$, or by $(\pi,\rho)$ where $\rho$ is a pentagonal
frame, have trivial stabilizer.  This verifies every single-orbit
row of (B.17).  The three choices of gommage at a fully pinned
$\mathrm {III}_1$ vertex give three arrows at each of its 120
members.  An equivariant orbit of arrows over a free transitive
source set contains exactly one arrow over every source: a
stabilizer of an arrow is contained in the trivial source
stabilizer.  Therefore these $360$ arrows split into precisely
three free orbits.  Adding gives the claimed sixteen orbits.

Two scan checks are essential here.  The long bent arrow labeled
$\gamma_r^2$ starts at $\mathrm {III}'_2$, not at
$\mathrm {III}_1$ as in the supplied TikZ transcription.  Its
end-valencies are $2,2$, agreeing with $120=2\cdot60=2\cdot60$
on scan 33.  The diagonal $\gamma_r^{1'}$ on scan 31 has valency
three at both of its 120-member ends.  Scan 33 retains
$3\cdot120$ on its left but adds the arithmetically incompatible
$3\cdot40=3\cdot40$ on its right.  Both the valencies and the
retained subtotal $1020$ force $360$ for that row.  We record the
incompatible reading rather than silently replacing its $40$.

The program \texttt{b\_finite\_checks.py} enumerates all $120$
permutations and all the displayed cosets.  It produces each
vertex and each member of each named arrow family, independently
reproducing (B.16)--(B.17).  Its endpoint maps are explicitly typed
coset maps $gH\mapsto(gH_s,gkH_t)$ with
$H\subset H_s\cap kH_tk^{-1}$.  The generated
\texttt{b\_diagram.json} states its precise scope: it is a
realization of these endpoint valencies, and does not assign the
geometric homotopy classes of the manuscript's paths.
\end{proof}

\begin{result}{B.13}{Proved}{The relations required by the five-point diagram}
The underlying geometric graph of the diagram in (B.16)--(B.17)
is connected.  The free groupoid on its displayed arrows has
isotropy group $F_{1079}$.  The fundamental group of $M_{0,5}$ has
abelianization $\mathbb Z^5$.  Therefore a presentation of
$\Pi_1M_{0,5}$ on exactly those vertices and arrows must impose
relations whose images in the abelianization of $F_{1079}$ span
a subgroup of rank $1074$.
\end{result}
\begin{proof}
The bipartite subgraph between $\mathrm I_6$ and $\mathrm I_4$
has a concrete coset description.  Label $K=\{0,1,2,3,4\}$ and take
\[
 H_6=\langle(012),(12)(34)\rangle,
 \qquad H_4=\langle(1324)\rangle,
 \qquad H_2=\langle(12)(34)\rangle.
\]
Its edges are $S_5/H_2$, with endpoints $gH_6,gH_4$.
The distinguished square vertex is $0$, and its diagonals are
$\{1,2\},\{3,4\}$; thus this is exactly the incidence obtained by
choosing the distinguished member of the three-element equator
of a bitetrahedron.  The group $\langle H_6,H_4\rangle$ is $S_5$.
For an elementary verification, conjugating $x=(012)$ by
$(1324)$ gives $y=(034)$.  Conjugating $x$ by $y,y^2$ gives
$(312),(412)$.  Thus all $(i12)$, $i\in\{0,3,4\}$ occur,
and products $(i12)(j12)^{-1}=(i1j)$ give all three-cycles
containing $1$.  The identity $(1ab)(1bc)=(abc)$ then gives
every three-cycle, hence $A_5$.  The four-cycle is odd.  Enumeration in the attached program also
checks the order $120$.  A path in this bipartite coset graph is a
product of elements of $H_6,H_4$, proving connectedness.  Every
other vertex type has a displayed arrow toward this subgraph or
toward a type already connected to it.  In particular
$\mathrm I_{10}$ connects through $\mathrm {II}_2$ and its
$\tau$-bijection with $\mathrm {II}'_2$.  Each incidence has
positive valency at every endpoint, so this proves connectedness
of the full graph.

Choose a spanning tree of the graph.  Successively contracting
its 541 edges gives a groupoid equivalence to the one-object free
group on $1620-541=1079$ generators, one for each remaining edge.
The comparison sends an arrow $e:v\to w$ to the loop obtained by
following the tree path from a fixed base vertex to $v$, then $e$,
then the inverse tree path from the base vertex to $w$.  This also
specifies the equivalence on every arrow.

Normalize three of the five markings to $0,1,\infty$.  Then
\[
 M_{0,5}=\{(x,y)\in\mathbb C^2:
 x,y\ne0,1,\quad x\ne y\}.
\]
Forgetting $y$ is a locally trivial bundle with base
$\mathbb C\setminus\{0,1\}$ and fibre
$\mathbb C\setminus\{0,1,x\}$.  A section is
$y=2+|x|$, since it avoids $0,1,x$.  It gives
$\pi_1(M_{0,5})=F_3\rtimes F_2$.  The base pure braids act trivially
on the abelianization of the fibre: each labelled puncture
meridian is sent to a conjugate of itself.  Abelianizing this
semidirect product therefore gives $\mathbb Z^3\oplus\mathbb Z^2$.
The rank statement follows from the exact abelianization of a
group presentation.  It requires at least $1074$ independent
abelianized relators; it does not purport to list a complete
nonabelian relation system.
\end{proof}

\begin{result}{B.14}{Proved}{The pentagon as an explicit relation among pants moves}
For a five-holed sphere with a fixed cyclic order and a fixed
outer boundary, let $T_0,\ldots,T_4$ be the five binary planar
pants trees.  Choose the standard compatible marking on each.
Write $A_i:T_i\to T_{i+1}$ for the elementary change of one
internal cut, with indices modulo five.  In the marking groupoid
of \citet[Sections 4.5, 4.7 and 4.10]{B_BakalovKirillov}, each
$A_i$ is the specified composite of a cut-removal move and a
cut-insertion move, and
\[
 A_4A_3A_2A_1A_0=1_{T_0}. \tag{B.18}
\]
Under the equivariant comparison of B.7 this is a genuine
relation in the mapping-class-group groupoid.  Under capping
boundary circles it remains a relation in $\Pi_1M_{0,5}$.
\end{result}
\begin{proof}
Let $U$ be the uncut marking with the chosen cyclic order and
outer distinguished edge.  For each $T_i$, remove its two cuts
and denote the resulting path to $U$ by $p_i$.  The generalized
$F$-move is defined in the cited Section 4.5 with a specified
distinguished edge at its target; hence all these maps have
specified objects and target marking.  The associativity-of-cuts
relation, equation (4.7) of that paper, says that the two orders
of removal give the same $p_i$.

Adjacent trees share one cut.  Let $N_i$ be the one-cut marking
obtained from either by removing the nonshared cut, and write
$f_i:T_i\to N_i$ and $f'_i:T_{i+1}\to N_i$ for the respective
cut-removal moves.  Set $A_i=(f'_i)^{-1}f_i$.  Removing the remaining
cut from $N_i$ gives a path $b_i:N_i\to U$.  Associativity gives
$p_i=b_if_i$ and $p_{i+1}=b_if'_i$, and hence
\[
 p_{i+1}A_i=p_i,
 \qquad A_i=p_{i+1}^{-1}p_i.
\]
Multiplying these five equations cancels the adjacent $p_i$'s
and proves (B.18).  Every move in $A_i$ is supported on the
four-holed subsurface where the pants decomposition changes.
The relation uses the common uncut five-holed marking; it is the
specific dimension-two compatibility absent from a presentation
consisting only of unrelated one-dimensional pieces.

More explicitly, the equivariant groupoid has arrows $(g,[p])$
with $p:gM\to M'$.  The word in (B.18) maps to
$(1,[A_4\cdots A_0])=(1,[1])$ by the homotopy just written.
Applying the boundary-twist quotient functor in B.7 preserves
this equality.  Thus both comparison maps and the path homotopy
are specified; no identification of an elementary move with an
unbased homeomorphism is used.
\end{proof}

\begin{result}{B.15}{Open}{The remaining literal path and sphere-only comparisons}
B.25 completes a specified geometric realization with all
$542$ objects and $1620$ arrows.  Its $1080$ defining relators
are complete, all arrow words and tree transports are computed,
and its relative-dimension-one dictionary is explicit.
B.26 proves that this relator count is minimal for the graph.
The exact covering functor in B.23 and the retained-circle
constructions in B.24--B.25 state which operations are used.

What remains is the identification of these specified conventions
with any further convention intended by the original finite
marked diagram: the direct $f_{\vec Q,a}$ path, the retained-pin
erasure, and the three tetrahedral pin choices in the
$\gamma_r^{1'}$ family.  The supplied scans contain the stated
incidences and fractional twists; those data alone do not
select a relative homotopy class for every such choice.
The discrepancy $3\cdot120$ versus the added $3\cdot40$ is
retained in B.12 and the reading register.  The actual convention
maps of B.22--B.25, rather than the arbitrary endpoint realization
in \texttt{b\_diagram.json}, define the proved full model.

For the global assertion, B.27 proves the finite geometric
colimit with all five elementary types and its tree index in
genus zero.  The further reduction to the original sphere-only
geometric dictionary, with elliptic cut-changing and framed
coherence maps, remains the issue in B.8.  These are now the
precise unproved comparisons; no unassigned arrow word remains
inside the specified full realization of B.25.
\end{result}

\subsection{A direct reconstruction of the early four-point model}

\begin{result}{B.16}{Proved}{Thirty objects, forty-eight arrows, and the four-point word reduction}
The arrangement model underlying cote 146, scans 20, 27, 31 and
33 (folios 19, 26, 30 and 32), admits the following complete typed
presentation.  It has thirty objects, forty-eight essential
arrows and eighteen defining relations.  The expanded generator
table of scan 27, including its full-turn and half-turn auxiliary
arrows, has ninety-six arrows.  Its isotropy group is
\[
 \langle A,F,t\mid[t,A]=[t,F]=1\rangle=F(A,F)\times\mathbb Z(t).
 \tag{B.19}
\]
The central $t$ is one positive complex-scalar turn in the
arrangement complement.  Quotienting by it gives the four-point
moduli group $F_2$.  The explicit words below give a comparison
functor on every generator, including the auxiliary ones.
\end{result}

\begin{proof}
Here are the data without uncertain subscripts.  Indices $k$ lie
in $\mathbb Z/6$, and $\epsilon$ is $+1$ or $-1$.  In a Euclidean
plane $V$ place three lines at angles $0,\pi/3,2\pi/3$ and write
$J$ for its positive quarter-turn.  Put
\[
 X=V_{\mathbb C}\setminus\bigcup_{j=0}^2(L_j)_{\mathbb C},
 \qquad Q_k=(\cos((k+\tfrac12)\pi/3),
                   \sin((k+\tfrac12)\pi/3)).
\]
Choose $0<\eta<\pi/12$.  In the chamber containing $Q_k$, choose
$R_{k,+}$ at angle $(k+1)\pi/3-\eta$, and $R_{k,-}$ at angle
$k\pi/3+\eta$, on the unit circle.  Finally put
\[
 P_{k,\epsilon}=Q_k-i\epsilon JQ_k.
\]
There are six $Q$'s, twelve $R$'s and twelve $P$'s.  On each of
the two complex eigenlines of $J$, the six $P$'s are consecutive
scalar sixth-roots, exactly the pole circles and the three object
families of the original diagrams.  Changing the sufficiently
small $\eta$ supplies the canonical path transport within the
same chamber and does not change this presentation.

The four families of arrows, each with twelve members, are
\[
\begin{array}{c|c}
 b_{k,\epsilon}&Q_k\to R_{k,\epsilon}\\
 c_{k,\epsilon}&P_{k,\epsilon}\to Q_k\\
 a_{k,\epsilon}&R_{k,\epsilon}\to R_{k+\epsilon,-\epsilon}\\
 d_{k,\epsilon}&P_{k,\epsilon}\to P_{k+\epsilon,\epsilon}.
\end{array}
\tag{B.20}
\]
Take $b$ to be the straight segment in the real chamber, take
$c(s)=Q_k-i\epsilon(1-s)JQ_k$, and take
$d(s)=e^{\pi i s/3}P_{k,\epsilon}$.  The identity
$J P_{k,\epsilon}=i\epsilon P_{k,\epsilon}$ gives precisely the
stated endpoint of $d$.  For $a$, let $u,v$ be its source and target
and set
\[
 a(s)=\tfrac12(u+v)+e^{\pi i s}\tfrac12(u-v).
\]
This is a positive half-turn in the complex normal coordinate to
the intervening line.  None of these paths meets an excluded
complex line: for $b,c$ the appropriate nonzero real linear
functional has constant sign in its real part; for $d$ its real
part rotates through two adjacent chambers and its imaginary
part prevents a zero at their common wall.  For $a$ the two other
line functionals have modulus of their constant term strictly
larger than that of their rotating term, by the bound on $\eta$.

To match the manuscript's arrow orientations, let
$r(k,\epsilon)$ denote the incidence of the intervening wall with
the chamber containing $Q_k$.  Then $b,c,d$ in (B.20) are its
$\widetilde b,\widetilde c,\widetilde d$ with these incidences,
whereas our positive crossing is
\[
 a_{k,\epsilon}
   =\widetilde a_{r(k+\epsilon,-\epsilon)}^{-1}.
 \tag{B.21}
\]
This orientation is forced by its scan 31 mixed equation
$\widetilde a\,\widetilde b\,\widetilde c\,\widetilde d
 =\widetilde b\,\widetilde c$: the $\widetilde a$ there traverses
back from the terminal chamber of $\widetilde d$.  Thus we specify
an inverse and its changed source, rather than identify two
oppositely directed arrows.

The twelve mixed relations are
\[
 c_{k+\epsilon,\epsilon}d_{k,\epsilon}
 =b_{k+\epsilon,-\epsilon}^{-1}a_{k,\epsilon}
       b_{k,\epsilon}c_{k,\epsilon}.                 \tag{B.22}
\]
Define the paths from $Q_k$ to $Q_{k+3}$
\[
 H_{k,\epsilon}=
 c_{k+3,\epsilon}
 d_{k+2\epsilon,\epsilon}d_{k+\epsilon,\epsilon}d_{k,\epsilon}
 c_{k,\epsilon}^{-1}.
\]
The six further relations are
\[
 H_{k,+}=H_{k,-}\quad(k\in\mathbb Z/6).             \tag{B.23}
\]
They are the typed half-turn equality represented by the
manuscript's uncertain products of $\widetilde A$'s in scans
31 and 33.  Three consecutive sixth-turns give a half-turn;
six give a full turn.  These path lengths determine the invariant
meaning without assigning a new paleographic reading to an
uncertain exponent.

We first verify these equations geometrically.  Both sides of
(B.22), after removing their common initial $c$, lie in the domain
where the real parts of the two noncrossed line functionals have
their prescribed signs.  In coordinates normal and tangent to
the crossed wall this domain has the form
\[
 \{(z,w)\in\mathbb C^*\times\mathbb C:
          \operatorname{Re}w>c\,|\operatorname{Re}z|\}
\]
after a real linear rescaling, for some $c>0$.  Its projection to
$\mathbb C^*$ has convex fibres, so relative path classes are
classified by winding in $z$.  Both paths have the same endpoints
and a positive argument change of $\pi$ in $z$.  This proves
(B.22).  For (B.23), the square
\[
 (u,s)\longmapsto e^{\pi i s}(Q_k-iu\epsilon JQ_k),
 \qquad 0\leq u,s\leq1,
\]
lies in $X$: every excluded real line functional has a nonzero
real part before the scalar factor.  Its boundary identifies
$H_{k,\epsilon}$ with the scalar half-turn
$\mu_{Q_k}(s)=e^{\pi i s}Q_k$.  Both signs give the same path
class, as required.

These are all defining relations.  Eliminate the twelve $R$
objects using the arrows $b$, and the twelve $P$ objects using
$c$.  Define the resulting wall crossings
\[
 e_{k,\epsilon}
   =b_{k+\epsilon,-\epsilon}^{-1}a_{k,\epsilon}b_{k,\epsilon}.
\]
The mixed equations eliminate all twelve $d$'s.  What remains has
six chamber objects, twelve directed wall crossings, and the six
relations equating the two length-three galleries to the
opposite chamber.  This is the $A_2$ chamber presentation.  One
may check completeness directly by identifying the three lines
with $x_i=x_j$ in $\{x_1+x_2+x_3=0\}$: adding the centroid gives
an isomorphism $\mathbb C\times X\cong\operatorname{Conf}_3
(\mathbb C)$.  The six chamber objects are its six real orderings,
and its adjacent normal half-turns are the standard adjacent
braid crossings.  Generic braid diagrams generate all paths;
a generic homotopy between diagrams has only cancellation of
opposite crossings or the three-strand move
$\sigma_1\sigma_2\sigma_1=\sigma_2\sigma_1\sigma_2$.
Lifting that relation to each ordering gives exactly (B.23).
Equivalently these are the six vertices, twelve edges and six
two-cells of the arrangement complex described in
\citet[Section 4.1]{B_SalvettiSettepanella}.  This verifies that
the path functor from (B.20)--(B.23) to $\Pi_1(X)$ is fully
faithful and essentially surjective.

For a completely algebraic word check, put
$x_k=e_{k,+}$ and $y_k=e_{k,-}$, and contract the spanning tree
$x_0,\ldots,x_4$.  Set $x_5=t$, $y_0=A$, $y_5=F$.  The six
remaining relations solve as follows:
\[
\begin{array}{c|rrrrrr}
k&0&1&2&3&4&5\\ \hline
x_k&1&1&1&1&1&t\\
y_k&A&A^{-1}F^{-1}&F&At&A^{-1}F^{-1}&F.
\end{array}                                                   \tag{B.24}
\]
They leave precisely $[t,A]=[t,F]=1$.  For example, the first
three equations give $y_4=y_1=A^{-1}F^{-1}$ and $y_2=F$; the
fourth gives $y_3=At$; the last two impose the displayed
commutators.  Conversely substitution into all six equations
verifies them, so no relator has been lost.  Sending every $b,c$
to $1$ and every $a_{k,\epsilon},d_{k,\epsilon}$ to the value
of $e_{k,\epsilon}$ in (B.24) gives the announced functor to
(B.19).  Every eliminated arrow belongs to a specified spanning
tree or is solved by a displayed relation, proving that this
functor is an equivalence.

The scalar full turn is the product of six successive $d$'s,
and maps to $t$ for either sign.  Every $\lambda$ in scan 27
therefore maps to $t$.  The $\mu_{Q_k}$ maps to $1$ for
$k=0,1,2$ and to $t$ for $k=3,4,5$; its target is $Q_{k+3}$,
so this does not identify a half-turn with an identity on its
original object.  The $\mu_R$ words are obtained by transport
along the specified $b$ arrows.  This accounts for all thirty
$\lambda$'s and all eighteen $\mu_Q,\mu_R$'s in that table.
The $\lambda$-transport relations follow from centrality of $t$.

Finally the affine normalization
\[
 (z_1,z_2,z_3)\longmapsto
 \bigl(z_1,z_2-z_1,(z_3-z_1)/(z_2-z_1)\bigr)
\]
identifies the ordered configuration space with
$\mathbb C\times\mathbb C^*\times
 (\mathbb P^1\setminus\{0,1,\infty\})$.
A scalar full turn is the positive generator of its
$\mathbb C^*$ factor, and the projective quotient kills exactly
$t$.  This proves the claimed comparison with $M_{0,4}$.
The program \texttt{b\_four\_point.py} lists all objects,
all forty-eight arrows and all eighteen typed relations, checks
both endpoints and the words (B.24), and checks all twelve pole
full turns.  Its output also verifies the expanded count $96$.
\end{proof}

\begin{result}{B.17}{Proved}{Four leg framings and the dual scalar convention}
The free-sign induction of the thirty-object model B.16 from
$(\mathbb Z,\mu_2)$ to $(\mathbb Z^4,\mu_2^4)$ along the diagonal
has $240$ objects.  A product presentation with three additional
coordinate half-turn circles has $1104$ essential generators,
or $1488$ when retaining the auxiliary generators of scan 27.
Its automorphism group is $\mathbb Z^4\times F_2$, and its
quotient by all four leg framings is $\Pi_1M_{0,4}$.
These counts name this particular induced version of the early
model; finite equivalent presentations need not have the same
object count.

The two scalar conventions in cote 146, scans 18 and 50--53,
are related by an explicit dualization functor.  On the full
winding lattice it is multiplication by $-1$, and on the finite
sign group it is the identity.  Hence the positive scalar $t$
of B.16 is the inverse of the positive principal-bundle turn
when the manuscript calls the scalar action ``opposite'' on
scan 53.  The corresponding lattice pushout is respectively
$t\mapsto(1,1,1,1)$ or $t\mapsto(-1,-1,-1,-1)$, with this
comparison map specified.
\end{result}
\begin{proof}
The diagonal lattice is primitive.  Complete its generator
$w=(1,1,1,1)$ to the basis $w,e_2,e_3,e_4$ of $\mathbb Z^4$.
Modulo two, the same basis splits the sign extension into the
old antipodal action and three additional independent sign
coordinates.  Thus the induced finite groupoid is the product
of B.16 with three copies of the two-object half-turn circle
of B.2.  The basis is a stated choice, not a permutation-invariant
splitting.  This gives $30\cdot2^3=240$ objects and
\[
 48\cdot8+30\cdot(3\cdot8)=1104
\]
generators.  Using the unreduced $96$ in place of $48$ gives
$1488$.  The product relations consist of the lifted eighteen
relations, the commuting squares in the three circle coordinates,
and the squares commuting an old arrow with a new coordinate
half-turn.  The product universal property proves completeness.
Its central lattice is the pushout of $\mathbb Z(t)$ along the
stated diagonal, so its group is $\mathbb Z^4\times F_2$.
Killing that lattice and identifying its sign lifts gives the
projective quotient in B.16, hence $\Pi_1M_{0,4}$.

For the convention comparison, let $L$ be a complex line bundle.
For each nonzero vector $v$ there is exactly one complex-linear
functional $\ell_v$ with $\ell_v(v)=1$.  Thus
$v\mapsto\ell_v$ is a canonical continuous map
$L^{\times}\to(L^{\vee})^{\times}$.  It satisfies
$\ell_{\alpha v}=\alpha^{-1}\ell_v$, descends to the ray circles,
and sends a positive scalar loop to a negative scalar loop.
It commutes with multiplication by $-1$ and is natural under
linear isomorphisms of line bundles.  This defines the functor
on paths as well as on objects, proves its lattice and sign
formulas, and tracks the reversal that scan 53 states.
Applying this functor before the diagonal pushout gives precisely
the second displayed lattice map.  No unproved identification
between a line and its dual, with unchanged positive inertia,
is made.
\end{proof}

\subsection{A finite chamber model and its exact braid dictionary}

\begin{result}{B.18}{Proved}{Finite chamber arrows reduced to pure moduli generators}
Let $m\geq3$.  Define a finite groupoid presentation $\mathcal C_m$
with objects the $m!$ orderings $w$ of $\{1,\ldots,m\}$ and arrows
$e_{w,i}:w\to ws_i$, $1\leq i<m$, where $s_i$ interchanges the
entries in positions $i,i+1$.  At each object impose the lifted
commutation relation for $|i-j|>1$ and the lifted length-three
braid relation for $j=i+1$.  Also kill at the identity ordering
the loop
\[
 z_m=(\sigma_1\cdots\sigma_{m-1})^m.                 \tag{B.25}
\]
Then $\mathcal C_m\simeq\Pi_1M_{0,m+1}$.  The comparison is induced
by the actual adjacent half-turn paths in the ordered configuration
space, followed by affine normalization.  Its finite counts are
\[
\begin{array}{c|r|r|r|r}
 \text{space}&\text{objects}&\text{arrows}&
 \text{rank-two relations}&\text{central relation}\\\hline
 M_{0,4}&6&12&6&1\\
 M_{0,5}&24&72&72&1\\
 M_{0,6}&120&480&720&1.
\end{array}                                                   \tag{B.26}
\]
Every generator has an explicit image in the generators
\[
 A_{ij}=\sigma_{j-1}\cdots\sigma_{i+1}\sigma_i^2
             \sigma_{i+1}^{-1}\cdots\sigma_{j-1}^{-1},
 \qquad 1\leq i<j\leq m,\quad (i,j)\ne(1,2).       \tag{B.27}
\]
The algorithm below and \texttt{b\_genus\_zero.json} specify all
these words for $m=3,4,5$, and give the inverse on every generator
by (B.27).  These are the counts of this specified comparison
model; the $542$ objects of cote 148 retain their count in B.12.
\end{result}
\begin{proof}
Place $m$ distinct labelled points on the real axis.  The positive
normal half-turn interchanging two adjacent positions gives
$e_{w,i}$.  The configuration remains labelled: its endpoint has
the two labels in exchanged positions.  The generic-braid-diagram
argument of B.16 works with any number of strands.  The only
codimension-two changes are two independent crossings in either
order, and the three-strand move.  Consequently the groupoid
before (B.25) is the fundamental groupoid of
$\operatorname{Conf}_m(\mathbb C)$ on these $m!$ configurations.
This is also the two-skeleton of the type $A_{m-1}$ complex in
\citet[Section 4.1]{B_SalvettiSettepanella}.  There are
$(m-1)m!$ arrows and $\binom{m-1}{2}m!$ displayed rank-two
relations.  Some relations can subsequently be eliminated; the
table counts the given presentation.

The map
\[
 (z_1,\ldots,z_m)\longmapsto
 \left(z_1,z_2-z_1,
  \left(\frac{z_j-z_1}{z_2-z_1}\right)_{3\leq j\leq m}\right)
                                                               \tag{B.28}
\]
is an isomorphism of complex manifolds
\[
 \operatorname{Conf}_m(\mathbb C)
 \longrightarrow\mathbb C\times\mathbb C^*\times M_{0,m+1},
\]
where labels $1,2,m+1$ are normalized to $0,1,\infty$.
The loop $z_m$ rotates the whole configuration once; it generates
the fundamental group of the second factor.  Its conjugate at
any other object is the corresponding scalar turn, so the one
relation (B.25) kills the entire fibre lattice.  This proves the
geometric comparison and its exact kernel.

Here is a word algorithm specifying the functor without leaving
choices of paths implicit.  For an ordering $w$, let $T_w$ be the
positive permutation braid produced by moving the entry destined
for position $1$ left to that position, then the entry destined
for position $2$, and so on.  Thus $T_w$ is a specified path from
the identity ordering to $w$.  For $e=e_{w,i}$ with target $v$,
the based loop is
\[
 T_v^{-1}\,e\,T_w.                                  \tag{B.29}
\]
The program records its braid word in traversal order as
$T_w,\sigma_i,T_v^{-1}$.  More generally, for a pure braid $b$,
let $k=\operatorname{lk}_{12}(b)$, the signed crossing count of
strands $1,2$ divided by two.  Since
$\operatorname{lk}_{12}(z_m)=1$, replace $b$ by $z_m^{-k}b$.
This is the exact central correction for affine normalization.

Combing this corrected braid gives a word in (B.27).  To specify
combing, track the positions of the labelled strands while
reading the braid word and delete strand $m$.  Discard a crossing
involving that strand; decrement its index by one when the deleted
strand is to its left.  This gives a pure braid $b_0$ on $m-1$
strands.  Embed $b_0$ with a stationary final strand and put
$u=b_0^{-1}b$.  Then $u$ is in the kernel of deletion.  In the
Artin action on $F(X_1,\ldots,X_m)$ write
\[
 X_m^u=C X_m C^{-1},
\]
choosing $C$ by removing the inverse matching prefixes and
suffixes of the freely reduced word $X_m^u$ until only $X_m$
remains.  Delete every $X_m^{\pm1}$ from $C$ and freely reduce.
Replace $X_i^{\pm1}$ in the result by $A_{im}^{\pm1}$.  Recurse
on $b_0$.  This supplies the promised word in pure-band generators.

For completeness, the reason this algorithm returns the braid,
rather than just one invariant of it, is as follows.  Forgetting
the last point is locally trivial: disjoint small discs about
the preceding points and translations supported in these discs
give local trivializations.  Its fibre is the plane with $m-1$
points removed.  Induction, starting with $\mathbb C^*$, shows
that its base has zero $\pi_2$, so the fibre group injects.
The kernel is therefore the free group on the point-push loops
$A_{1m},\ldots,A_{m-1,m}$.  In the displayed Artin formula,
$A_{im}$ sends $X_m$ to $X_iX_mX_i^{-1}$.  On the kernel,
deleting $X_m$ from the conjugator is a homomorphism to
$F(X_1,\ldots,X_{m-1})$, since its action on that quotient free
group is the identity.  It takes the stated free basis to its
stated free basis.  It is consequently an isomorphism, which
proves the fibre step of combing.  Recursion proves the claim.
The final two-strand exponent is $\operatorname{lk}_{12}$;
it vanishes after the central correction, explaining the omission
of $A_{12}$ in (B.27).

This argument also proves directly that two outputs represent
the same mapping class exactly when their iterated free-group
normal forms agree.  The program makes an additional check of
each output in the faithful Artin action.  Faithfulness can be
seen by choosing disjoint arcs from a fixed boundary basepoint
to the punctures: an automorphism fixing the peripheral free
basis preserves the homotopy classes of these arcs, which can
be isotoped successively to the original arcs.  Cutting along
them leaves a disc, and the remaining boundary-fixing map is
isotopic to the identity.  Thus this check concerns the braid
itself.

Sending each $A_{ij}$ back to its path (B.27) is inverse to the
collapse functor on vertex groups.  Formula (B.29) and the paths
$T_w$ supply the natural isomorphism for the other objects.
Hence the maps specify an equivalence generator by generator.
All $12+72+480$ arrow outputs, the typed rank-two relations, and
their central corrections were computed and checked.  The
algorithm is also an explicit way to pass any additional pure
braid word, once assigned to a geometric sewing path, into this
finite presentation.
\end{proof}

\begin{result}{B.19}{Proved}{Exact comparison with the five- and six-point compact presentations}
Under (B.27)--(B.29), the five-point vertex group is
\[
 \Gamma_{0,5}=F(x,y,z)\rtimes F(a,b),
 \qquad (a,b,x,y,z)=(A_{13},A_{23},A_{14},A_{24},A_{34}).
\]
Put $c(u,v)=v^{-1}u^{-1}vu$.  Its six defining relations are
\[
\begin{array}{lll}
 axa^{-1}=z^{-1}xz,&
 aya^{-1}=c(x,z)y c(x,z)^{-1},&
 aza^{-1}=c(x,z)z,\\
 bxb^{-1}=x,&
 byb^{-1}=z^{-1}yz,&
 bzb^{-1}=c(y,z)z.
\end{array}                                                   \tag{B.30}
\]
For six points, adjoin
$(p,q,r,s)=(A_{15},A_{25},A_{35},A_{45})$.
Then $\Gamma_{0,6}=F(p,q,r,s)\rtimes\Gamma_{0,5}$, with the six
relations (B.30) and the twenty equations
$huh^{-1}=\phi_h(u)$ in the following table:
\begingroup\small\setlength{\arraycolsep}{3pt}
\[
\begin{array}{c|cccc}
 h&\phi_h(p)&\phi_h(q)&\phi_h(r)&\phi_h(s)\\\hline
 a&r^{-1}pr&c(p,r)q c(p,r)^{-1}&c(p,r)r&s\\
 b&p&r^{-1}qr&c(q,r)r&s\\
 x&s^{-1}ps&c(p,s)q c(p,s)^{-1}&c(p,s)r c(p,s)^{-1}&c(p,s)s\\
 y&p&s^{-1}qs&c(q,s)r c(q,s)^{-1}&c(q,s)s\\
 z&p&q&s^{-1}rs&c(r,s)s
\end{array}                                                   \tag{B.31}
\]
\endgroup
These are exactly the generators and relations in the supplied
EXIT\_ONE\_TYPE statements P5c2 and P5c3, with the comparison
specified by the individual band words (B.27).  Completeness
and the equations are independently proved here.  This comparison
does not assert that every orbit transporter for the separate
complex of groups on $P_2$ has already been expressed in these
words.
\end{result}
\begin{proof}
The normalization (B.28) identifies $M_{0,n}$ with the ordered
configuration space of $n-3$ points in
$\mathbb C\setminus\{0,1\}$.  Forget the last moving point.
The local trivializations just constructed prove the particular
fibration of \citet[Theorem 1]{B_FadellNeuwirth}.  A continuous
section places the last point at the positive real number
$2+\sum |z_i|$, beyond all previous points.  The base and the
punctured-plane fibre are aspherical by induction.  The exact
homotopy sequence therefore gives, with these basepoints and
section, the stated split extensions with free kernels of
ranks three and four.

The fibre generators are the indicated $A_{ij}$, as in the
combing proof.  Conjugate each of these band words by the band
word of a base generator and apply the fibre-conjugator
calculation of B.18.  It gives precisely (B.30) and (B.31).
The program checks the six equations in $B_4$ and the twenty new
equations in $B_5$ by their free-group actions.  It also constructs
the inverse fibre automorphisms by conjugating with inverse base
braids, and checks both inverse identities.  For each of the six
base relators and each of the four new fibre generators it
verifies that the induced automorphism is the identity as a
freely reduced word: twenty-four checks, with no abelianization.

Finally, a split extension of a free group by a presented group
has a presentation consisting of the base presentation and these
conjugation equations.  Indeed, use them to move base letters to
one side of any word.  Its image in the base first determines the
base factor; an element mapping to the identity is then a word
in the injective free kernel, so it vanishes exactly when that
free word reduces to the empty word.  This proves completeness
of the six and twenty-six relations.  The same normal form proves
that the inverse band-word maps of B.18 agree with these
presentations on every generator.  It supplies an independently
verified comparison, rather than relying on the status labels of
the previous return.
\end{proof}

\subsection{The involution curve and the sixty explicit core paths}

\begin{result}{B.20}{Proved}{A geometric realization of the bipyramid--pyramid incidence}
Let $U=\mathbb A^1_{\mathbb C}\setminus\{2,-2\}$ and
\[
 q_\mu(z)=z^4-\mu z^2+1.
\]
The roots of $q_\mu$, together with the distinguished point
$\infty$, define a family of five-pointed spheres with the
involution $z\mapsto-z$.  After adjoining an ordering of the
four roots, it has the explicit comparison to $M_{0,5}$ given
by (B.28).  The parameter $\mu=0$ belongs to this family:
its fibre is the square-pyramidal configuration.  Thus the
extra deletion of $0$ in cote 148, scan 39 (folio 38), does not
follow from the moduli problem with the fourth-root frame
specified on scans 37--38 (folios 36--37).

There is the following explicit realization in $\Pi_1M_{0,5}$
of the sixty $f_{B,s}$ incidences of scan 31.  In one labelling,
put
\[
 K(r)=\bigl(\infty,-ir,ir,-1,1\bigr),
 \qquad \frac1{\sqrt3}\leq r\leq1.                 \tag{B.32}
\]
Its initial configuration is bipyramidal and its final
configuration is square-pyramidal.  Relabelling (B.32) gives
sixty paths, with twenty initial objects, thirty final objects,
three paths leaving each initial object and two arriving at
each final object.  The endpoint functors are exactly
\[
 \begin{gathered}
 \mathfrak S_5/H_B \longleftarrow
 \mathfrak S_5/H \longrightarrow \mathfrak S_5/H_Q,\\
 H_B=\langle(012),(12)(34)\rangle,\quad
 H_Q=\langle(1324)\rangle,\quad
 H=\langle(12)(34)\rangle.
 \end{gathered}                                    \tag{B.33}
\]
In particular these are the endpoint maps used in B.13,
now furnished with actual paths.  Formula (B.32) specifies
one realization of this part of the manuscript diagram;
it does not specify the remaining $1560$ arrows or assert
an unrecorded homotopy convention for them.
\end{result}
\begin{proof}
An involution of a complex projective line has two fixed
points.  Choose the marked one as $\infty$ and the other as
$0$.  The involution becomes $z\mapsto-z$, and its two
unmarked orbits have squared coordinates $\xi,\eta\ne0$,
with $\xi\ne\eta$.  The fourth-power frame in the manuscript
normalizes $\xi\eta=1$.  Writing $\mu=\xi+\eta$ gives
$q_\mu$ and the single condition $\mu^2\ne4$.  Conversely
the roots of this polynomial recover both orbits and their
frame.  Its discriminant is
\[
 \operatorname{disc}_z(q_\mu)=16(\mu^2-4)^2.         \tag{B.34}
\]
Changing the frame by $\zeta\in\mu_4$ changes $\mu$ by
$\zeta^2\mu$.  In particular $\mu=0$ is a valid framed
object.  Its roots are a scalar multiple of $\mu_4$; it
has the extra order-four symmetry, but this does not make
it singular or invalidate its specified order-two symmetry.
The finite cover
\[
 \{(\mu,z_1,z_2,z_3,z_4):
     q_\mu(z_j)=0,\ z_j\ne z_k\ (j\ne k)\}
 \longrightarrow U
\]
retains all labellings.  Sending its points to the labelled
sphere $(z_1,z_2,z_3,z_4,\infty)$, or to its normalized
coordinates (B.28), defines the promised map and the induced
fundamental-groupoid functor.

For the incidence, label the standard bipyramid
\[
 (1,\zeta,\zeta^2,0,\infty),\qquad
 \zeta=\exp(2\pi i/3),
\]
by $0,1,2,3,4$ in that order, and select its equatorial
point $s=0$ (the coordinate $1$).  The transformation
$M(z)=(z+1)/(z-1)$ takes it to
$K(1/\sqrt3)$.  Along (B.32) no points collide.  At the
other endpoint the four points in the finite plane are
$\{1,i,-1,-i\}$, with their complex cyclic order; the
selected point is the apex.  The bipyramid's group acts
as $H_B$ on the five labels.  At the final endpoint the
positive quarter rotation acts by $(4\,2\,3\,1)$,
which generates the same subgroup as $(1324)$.
The common half rotation is $(12)(34)$.

That common involution preserves the whole path.  No other
permutation can identify two relabelled paths: it would
preserve both endpoints, and
$H_B\cap H_Q=H$.  Conversely every element of $H$ does
identify them, by $z\mapsto-z$.  Hence the paths are
indexed by $\mathfrak S_5/H$, with endpoints the cosets
in (B.33).  This gives their cardinalities and valencies
without choices of unidentified conjugating paths.

To see the relation with the manuscript's quartic
coordinate, rescale the finite-plane coordinate in (B.32)
so that the product of the two squared orbits is $1$.
One obtains
\[
 \mu=i(r-r^{-1}),\qquad
 \mu(1/\sqrt3)=-2i/\sqrt3,\qquad \mu(1)=0.
\]
This is precisely a segment of the family just constructed.
The program \texttt{b\_quartic.py} verifies the fractional
linear transformation, discriminant, three stabilizers,
all sixty endpoint maps, and both valencies over exact
algebraic numbers.  Thus the comparison on this subdiagram
is geometric and fully specified.  Matching all the other
named arrows and their relations remains the separate
calculation recorded in B.15.
\end{proof}

\begin{result}{B.21}{Proved}{A complete presentation on the sixty geometric core paths}
The fifty-object, sixty-arrow path diagram (B.32)--(B.33)
has a quotient groupoid equivalent to $\Pi_1M_{0,5}$ by the
following twelve relations.  Its comparison functor sends
every arrow to the path (B.32), with its specified relabelling.
The construction below specifies all basepoint transports,
eleven free graph cycles, their five-point braid words, and
an inverse on the five standard generators.

Number cosets of each group in (B.33) by the lexicographically
least permutation in the coset, with the permutations written
as tuples of images of $0,1,2,3,4$.  Number the sixty arrows
$E_1,\ldots,E_{60}$ in that order.  Use the breadth-first
spanning tree from the first bipyramidal vertex, inspecting
incident arrows in their numerical order.  Its non-tree
arrows are
\[
 E_2,E_3,E_6,E_8,E_9,E_{13},E_{19},E_{22},E_{42},E_{45},E_{48}.
                                                               \tag{B.35}
\]
If $t_v$ is the tree path from the root to $v$, the corresponding
free cycles $g_j$ traverse $t_{s(E)}$, then $E$, then
$t_{t(E)}^{-1}$.  Their exact words are the following.  A signed
integer $k$ denotes $E_k$ for $k>0$ and $E_{-k}^{-1}$ for $k<0$;
within each bracket the paths are traversed from left to right.
\begingroup\small
\[
\begin{array}{c|l}
g_1 &[1,-12,41,-44,20,-17,50,-59,35,-26,2,-11,53,-56,33,-28,39,-46,23,-14]\\
g_2 &[1,-12,54,-55,21,-16,3,-10,52,-57,23,-14]\\
g_3 &[1,-12,41,-44,20,-17,6,-7,39,-46,23,-14]\\
g_4 &[1,-12,54,-55,34,-27,8,-5,49,-60,36,-25]\\
g_5 &[1,-12,41,-44,31,-30,9,-4,38,-47,36,-25]\\
g_6 &[1,-12,41,-44,20,-17,50,-59,35,-26,13,-24,58,-51,30,-31,44,-41,12,-1]\\
g_7 &[14,-23,57,-52,29,-32,19,-18,49,-60,36,-25]\\
g_8 &[14,-23,46,-39,28,-33,22,-15,38,-47,36,-25]\\
g_9 &[14,-23,46,-39,28,-33,56,-53,42,-43,32,-29,52,-57,23,-14]\\
g_{10}&[1,-12,41,-44,31,-30,51,-58,45,-40,27,-34,55,-54,12,-1]\\
g_{11}&[1,-12,41,-44,20,-17,50,-59,48,-37,16,-21,55,-54,12,-1].
\end{array}                                                    \tag{B.36}
\]
\endgroup
With the transport to the normalized configuration
$(0,1,10,20,\infty)$ specified in the proof, their images
in the presentation (B.30) are
\[
\begin{array}{c|l@{\qquad}c|l}
g_1&abyz&g_2&z\\
g_3&abz&g_4&z^{-1}y^{-1}z\\
g_5&axz&g_6&z^{-1}x^{-1}z\\
g_7&z^{-1}xz&g_8&b^{-1}z^{-1}y^{-1}\\
g_9&b^{-1}a^{-1}z^{-1}y^{-1}x^{-1}&g_{10}&abxyz\\
g_{11}&a&&.
\end{array}                                                    \tag{B.37}
\]
Define words in these cycles by
\[
 A=g_{11},\qquad B=g_{11}^{-1}g_3g_2^{-1},\qquad
 X=g_2g_6^{-1}g_2^{-1},\qquad
 Y=g_2g_4^{-1}g_2^{-1},\qquad Z=g_2.              \tag{B.38}
\]
The first six defining relations are
\[
\begin{array}{lll}
 g_1=g_3g_4^{-1},&g_5=g_{11}g_2g_6^{-1},&g_7=g_6^{-1},\\
 g_8=g_2g_3^{-1}g_{11}g_4g_2^{-1},&
 g_9=g_2g_3^{-1}g_4g_6g_2^{-1},&
 g_{10}=g_3g_6^{-1}g_4^{-1}.
\end{array}                                                    \tag{B.39}
\]
The remaining six, with $c(U,V)=V^{-1}U^{-1}VU$, are
\[
\begin{array}{lll}
 AXA^{-1}=Z^{-1}XZ,&
 AYA^{-1}=c(X,Z)Yc(X,Z)^{-1},&
 AZA^{-1}=c(X,Z)Z,\\
 BXB^{-1}=X,&
 BYB^{-1}=Z^{-1}YZ,&
 BZB^{-1}=c(Y,Z)Z.
\end{array}                                                    \tag{B.40}
\]
All twelve are based at the root; their conjugates supply
the relation at any other object.  Formulas (B.36),
(B.38)--(B.40) are a complete explicit system, with no
additional unspecified relators.
\end{result}
\begin{proof}
\emph{Actual configuration paths and transports.}
For an arrow coset choose its lexicographically least
representative $p$.  Relabel (B.32) by $p$, so that the
coordinate of label $j$ is $u_j(r)=K_{p^{-1}(j)}(r)$.
Let
\[
 N_{a,b,c}(w)=\frac{(w-a)(b-c)}{(w-c)(b-a)},
                                                               \tag{B.41}
\]
with its usual projective limits when one of the entries
is $\infty$.  The path in the normalized configuration space
is
\[
 Z_p(r)=\bigl(0,1,x_p(r),y_p(r)\bigr),\qquad
 (x_p,y_p)=\bigl(N_{u_0,u_1,u_4}(u_2),
                  N_{u_0,u_1,u_4}(u_3)\bigr).       \tag{B.42}
\]
All sixty rational functions and endpoint configurations
are determined by (B.41)--(B.42).  The endpoint configuration
depends only on the relevant $H_B$ or $H_Q$ coset, because
normalization is invariant under a common projective
transformation.  Thus these are the actual endpoint functors,
not arbitrary compatible maps of coset sets.

For each of the fifty endpoint configurations $(0,1,x,y)$,
choose the following path $T_{x,y}$ from $(0,1,10,20)$.
First move $10$ to $x$ through $N+iN^2$, leaving $0,1,20$
fixed.  Here $N\geq2$ is the least integer for which both
straight segments avoid $\{0,1,20\}$.  Then move $20$ to
$y$ through $N'+i(N')^2$, leaving $0,1,x$ fixed, with the
same least-integer rule.  Omit a motion whose endpoints
coincide.  None of these fifty values of $x$ is $20$, as
substitution in (B.42) verifies.  These choices always
exist: each forbidden collinearity puts a quadratic
condition on the integer $N$, with at most two roots;
a line cannot contain the parabola $N+iN^2$.  The endpoints
themselves avoid the stationary points.  This specifies
all transports without a homotopy-class choice left over.
The based image of an arrow, in traversal order, is
\[
 T_{s(E)}\,Z_p\,T_{t(E)}^{-1}.                     \tag{B.43}
\]
For a core cycle the intermediate transports cancel, leaving
the transport of that actual cycle to the global basepoint.

\emph{Certified straightening.}
Each rational path (B.42) is homotopic, relative endpoints,
to the straight segment between its endpoint configurations.
Here is a uniform proof of the homotopy, also checked by
exact polynomial certificates.  For any pair of its four
points, let $D(r)$ be their nonzero difference and put
$c=3/4$.  The difference formula for the fractional linear
map (B.41) expresses $D(r)$ as a quotient with at most three
original point differences in the numerator and three in
the denominator; limits at infinity only remove factors.
The original differences are constant real or imaginary
multiples, or multiples of $1+ir$ and $1-ir$.
On $1/2\leq r\leq1$ each nonconstant factor changes its
argument from its value at $c$ by at most
\[
 \max\{\arctan(3/4)-\arctan(1/2),
           \arctan(1)-\arctan(3/4)\}
   =\arctan(2/11)<\frac{\pi}{12}.
\]
For the last inequality, $2/11<2-\sqrt3=\tan(\pi/12)$;
squaring the positive quantities proves it.  Hence the
argument of $D(r)/D(c)$ has absolute value less than
$\pi/2$, and
\[
 \operatorname{Re}\bigl(D(r)\overline{D(c)}\bigr)>0.
                                                               \tag{B.44}
\]
The endpoint chord difference has the same strict
half-plane property, since it is a convex combination
of $D(1/\sqrt3)$ and $D(1)$.  Interpolating between the
rational path and its chord consequently never produces
a collision.  This proves the claimed homotopy.

The program \texttt{b\_core\_paths.py} also checks (B.44)
for all six pair differences on each of the sixty paths.
It writes the numerator and denominator of the real
rational function and uses the exact Sturm root count on
$[1/2,1]$ to certify their constant nonzero signs.  All
sixty paths require only one chord.  Thus the subsequent
crossing calculation represents the specified paths,
with a proved homotopy behind the replacement.

\emph{The braid computation.}
Project a complex point to
$\operatorname{Re}(z)+\operatorname{Im}(z)/37$.
All endpoints of the chords and transport segments have
distinct projected positions.  Their coordinates lie in
$\mathbb Q(i,\sqrt3)$.  On each linear segment, every
crossing time is obtained by solving a linear equation
over $\mathbb Q(\sqrt3)$.  Its order and sign are exact:
the sign of $a+b\sqrt3$ is decided from the signs of
$a,b$ and the comparison of $a^2$ with $3b^2$.
At a crossing the two strands are adjacent in the current
projected order.  In the positively oriented perpendicular
coordinate $\operatorname{Im}(z)-\operatorname{Re}(z)/37$,
the positive generator has the previously left strand
below the other strand.  This agrees with the positive
normal half-turn in B.18.  Simultaneous independent
crossings commute; no crossing involving three strands
occurs in the chosen projection.

This procedure gives each of the sixty words (B.43) in
$B_4$.  They are pure, and their linking number on the
stationary strands $0,1$ is zero.  Thus their images in
$\Gamma_{0,5}$ require no further scalar correction.
Combing by B.18 gives words of length at most four in
$a,b,x,y,z$; their braid-word lengths are at most ten.
The program verifies each combed word by the exact Artin
action.  It then computes the spanning tree and the eleven
cycles (B.36), and gives exactly the table (B.37).
The full arrow words, endpoint configurations, tree and
polynomial certificates are in
\texttt{b\_core\_paths.json}; the program is sufficiently
specified above to reproduce each entry.

\emph{Generation and completeness of the relators.}
Substitution of (B.37) in (B.38) gives exactly
$(a,b,x,y,z)$.  This proves surjectivity on the vertex
group and supplies an explicit inverse on its five
standard generators.  The graph has fifty vertices and
sixty edges and is connected, so its free vertex group
has the eleven displayed generators.  Each equation
(B.39) is immediate by substitution in (B.37).
Use these six equations to eliminate
$g_1,g_5,g_7,g_8,g_9,g_{10}$.  The five remaining free
generators are $g_2,g_3,g_4,g_6,g_{11}$.
The change of free basis (B.38) has the explicit inverse
\[
 g_2=Z,\qquad g_3=ABZ,\qquad
 g_4=Z^{-1}Y^{-1}Z,\qquad
 g_6=Z^{-1}X^{-1}Z,\qquad g_{11}=A.
\]
After that change the remaining relations (B.40) are
precisely the complete presentation (B.30), proved in
B.19.  Therefore the comparison on the vertex group is
an isomorphism.  All fifty objects lie in the same component,
and the tree paths give the natural isomorphisms at all
objects.  This proves the asserted equivalence of groupoids.

The independent program \texttt{b\_core\_generation.py}
checks every substitution and every relator in the exact
$F_3\rtimes F_2$ normal form.  As a further direct check
of generation, it folds the eleven base words in $F_2$
while retaining words in the original generators, obtains
the two base generators, and folds the resulting kernel
words in $F_3$.  The latter graph is already the
three-petal rose, so no saturation step is required.
It records these witnesses as well as the shorter inverse
(B.38).  The resulting twelve nonempty relators and their
expanded words in the sixty original edges are retained
in \texttt{b\_core\_generation.json}.  The proof of their
completeness is the explicit elimination just given.
\end{proof}

\begin{result}{B.22}{Proved}{A geometric extension through the square--triangle families}
There is a specified realization of the five object families
\[
 \mathrm I_{10},\quad\mathrm I_6,\quad\mathrm I_4,
 \quad\mathrm {II}_2,\quad\mathrm {II}'_2
\]
and the five arrow families
\[
 f_{B,s},\quad f_{\vec Q,a},\quad\tau_{\pi,s},
 \quad\gamma_{\pi,s},\quad\gamma''_{\pi,s}
\]
with the following
counts and complete presentation:
\[
\begin{array}{c|r}
 \text{objects}&12+20+30+60+60=182\\
 \text{arrows}&60+120+60+60+60=360\\
 \text{free cycles after a spanning tree}&179\\
 \text{defining relators}&180.
\end{array}                                                     \tag{B.45}
\]
The comparison sends these paths to $\Pi_1M_{0,5}$ and is an
equivalence after imposing the specified relators.  The
quarter-turn and the two erasures are furnished by the
one-node plumbing family below.  We specify $f_{\vec Q,a}$
as an erasure followed by a core path; this records its
path convention, in addition to its endpoint incidence.
All $360$ braid images and all $180$ relators are computed
in \texttt{b\_first\_tier.json}.
\end{result}
\begin{proof}
Use
\[
 V(q)=(\infty,-q,q,-1,1),\qquad
 q\in\mathbb C\setminus\{0,1,-1\},\qquad
 \epsilon=1/10.                                      \tag{B.46}
\]
At $q=0$ this has a stable limit consisting of the square
four-point component $(\infty,-1,1,0)$ and a three-point
tail.  Indeed, on the tail put $w=z/q$, so its marks are
$-1,1,\infty$, and the node coordinates $z,v=1/w$ satisfy
$zv=q$.  Thus $q$ is the actual normal plumbing coordinate.
For $|q|\leq\epsilon$ one may choose disjoint coordinate
discs of radius $1/2$ at the two branches.  The points
$V(\epsilon)$ and $V(i\epsilon)$ give smoothing
representatives of the right and left square--triangle
objects.  Both have the stabilizer $H$ in (B.33), so each
has sixty labelled realizations.

The twist is the positive quarter-turn
$q(t)=\epsilon\exp(\pi it/2)$; in the braid computation it
is replaced by
$\epsilon(1-t+it)$.  The radial interpolation between these
two paths avoids $0,1,-1$, proving the required endpoint
homotopy before normalization.  The two erasures are the
positive real and imaginary intervals
\[
\begin{split}
 \gamma_{\pi,s}:&\quad
  q:\epsilon\longrightarrow\sqrt5-2,\\
 \gamma''_{\pi,s}:&\quad
  q:i\epsilon\longrightarrow i/\sqrt3.
\end{split}                                                     \tag{B.47}
\]
They begin in the indicated normal ray, then continue
in the displayed smooth family.  The imaginary endpoint
is the bipyramid of B.20.  The real endpoint is pentagonal.
To check this last assertion, the transformation
$(z+1)/(z-1)$ sends the fifth roots of unity to $\infty$
and the four imaginary points of magnitudes
$\cot(\pi/5)$ and $\cot(2\pi/5)$.  Their ratio is
\[
 \frac{\cot(\pi/5)}{\cot(2\pi/5)}=2+\sqrt5.
\]
After scaling the larger magnitude to $1$ the smaller
is $\sqrt5-2$, giving (B.47).  Its cyclic order on the
real projective line is $(0,3,1,2,4)$, up to reversal,
and its stabilizer is
\[
 H_P=\langle(0\,3\,1\,2\,4),(12)(34)\rangle\simeq D_5.
\tag{B.48}
\]
Consequently the real erasure has the exact endpoint map
$\mathfrak S_5/H\to\mathfrak S_5/H_P$, with five incoming
arrows per pentagon.  The imaginary erasure has
$\mathfrak S_5/H\to\mathfrak S_5/H_B$, with three incoming
arrows per bipyramid.  The twist is the identity on the
common coset parameter $\mathfrak S_5/H$.

For clarity, the $120$ choices denoted $f_{\vec Q,a}$
are specified as follows.  A left square--triangle object
is the pair $(B,P)$, with $P$ one of the three equatorial
marks of its bipyramid after erasure.  Choose one of the
other two equatorial marks $s$, and use the composite
\[
 (B,P)\xrightarrow{\gamma''}B\xrightarrow{f_{B,s}}Q.
                                                               \tag{B.49}
\]
The apex of $Q$ is $s$.  The remaining equatorial pair
of $B$ is a diagonal of its oriented square.  Starting
at $P$ and following the cyclic orientation gives a side
$a$ of $Q$.  Conversely a side $a$ of an oriented square,
with initial vertex $P$, determines that diagonal, one
of the two incident bipyramids, and the pair $(B,P)$.
This is a bijection between (B.49) and the parameters
$(\vec Q,a)$: two choices per source and four per target.
It proves the finite endpoint comparison as well as
specifying the path used.  The source does not give a
separate parameterization of its direct $f_{\vec Q,a}$
path; (B.49) is the explicit convention of this realization.

Here is the complete relator prescription, including its
finite indexing and its computed word data.  Keep the first
fifty objects and sixty arrows of B.21.  Append the twelve
pentagons, sixty right square--triangles and sixty left
square--triangles, in lexicographic coset order.  Append
the sixty real erasures, sixty twists, sixty imaginary
erasures, then the $120$ composites (B.49), in that order.
The latter are ordered by their source index and then by
the core-edge index.  This numbers every object and arrow
without a choice.

Retain B.21's forty-nine-edge tree.  In increasing order,
add a new edge to the tree precisely when it joins two
previously different components.  The result has $181$
edges and all $182$ vertices.  Write $t_v$ for its unique
rooted path.  For each non-tree edge $E$, let
\[
 g_E=t_{s(E)}E t_{t(E)}^{-1}
\]
in traversal order, and compute its word $W_E(a,b,x,y,z)$
by the normalized-path and combing algorithms of B.21.
For the $168$ new non-tree edges impose exactly
\[
 g_E= W_E(A,B,X,Y,Z),                                \tag{B.50}
\]
where (B.38) gives the words on the right in the original
eleven core cycles.  Add the twelve core relators
(B.39)--(B.40).  These are the $180$ relators in (B.45).
The file \texttt{b\_first\_tier.json} lists every $E$,
$W_E$, tree path, and both sides of each expanded relation;
none of the $168$ words in (B.50) is left as an uncomputed
input.

For the path calculation, the normalized coordinates are
again (B.41)--(B.42), with $r=-iq$.  The real erasures
have no projected crossings, since all four finite
coordinates remain real and distinct.  For each pentagon
choose its transport through the first incoming real
erasure and its already specified right-object transport.
The right and left object transports use B.21's same
least-integer detour rule.  The sixty twists and sixty
imaginary erasures have rational coordinate paths;
all $120$ are homotopic to one chord by the strict
half-plane tests recorded in the program.  Exact Sturm
counts, exact crossing signs over $\mathbb Q(\sqrt3)$,
and Artin combing verify their words.  Pentagon coordinates
in $\mathbb Q(\sqrt5)$ are checked algebraically for the
shared endpoints; they introduce no crossings or numerical
approximations in the braid calculation.  Each composite
(B.49) is then checked by multiplying its two verified
words.

Finally, collapse the chosen tree.  The resulting free
group has $360-182+1=179$ generators: eleven old and
$168$ new.  Relations (B.50) eliminate all the new
ones.  The remaining presentation is exactly B.21,
which is $\Gamma_{0,5}$.  The comparison and its inverse
on vertex groups are therefore the actual path functor
and (B.38).  Tree transports extend the inverse to the
whole groupoid.  This proves both the relator prescription
and its completeness.
\end{proof}

\begin{result}{B.23}{Proved}{The realized paths come from an explicit four-point double-cover functor}
Put
\[
 U=\mathbb P^1_q\setminus\{0,1,-1,\infty\},\qquad
 p:U\longrightarrow M_{0,4},\quad p(q)=q^2,
\]
using the coordinate $M_{0,4}=\mathbb P^1_\lambda
\setminus\{0,1,\infty\}$.  There are exact functors
\[
 \Pi_1U\xrightarrow{p_*}\Pi_1M_{0,4},\qquad
 \Pi_1U\xrightarrow{D_*}\Pi_1M_{0,5},\quad
 D(q)=(\infty,-q,q,-1,1).                         \tag{B.51}
\]
The first is the two-sheeted covering functor.  The second
is induced by the double cover of the four-marked sphere
branched at $0,\infty$, retaining the point above $\infty$
and the four ordered points above $1,\lambda$.
For every arrow of the realized $360$-arrow model, its
image is a composite of explicitly relabelled images
under $D_*$ of lifted four-point paths.  Thus this is a
proved realization through a family of relative dimension
one, with the covering and forgetting maps retained.
It does not identify an unspecified sewing-only convention
for the other arrows in the margin of scan 33.

For a positive real basepoint $0<q<1$, order the four
finite points as $-1,-q,q,1$.  If $u_0,u_+,u_-$ are the
standard meridians at $q=0,1,-1$ specified below, their
braid images are
\[
 u_0\longmapsto\sigma_2^2,\qquad
 u_+\longmapsto\sigma_1^2\sigma_3^2,\qquad
 u_-\longmapsto\sigma_2\sigma_1^2\sigma_3^2\sigma_2^{-1}.
                                                               \tag{B.52}
\]
Under the scalar quotient and B.19's generators these are
\[
 b,\qquad b^{-1}a^{-1}z^{-1}y^{-1}x^{-1}z,
 \qquad az^{-1}yz,                                  \tag{B.53}
\]
respectively.
\end{result}
\begin{proof}
The map $q\mapsto q^2$ is unramified on the displayed
open curves and has two points above each $\lambda$.
The universal four-point sphere is normalized to
$(0,\infty,1,\lambda)$.  Its double cover $z^2=w$
has branch points $0,\infty$, and the four chosen lifts
are $\pm1,\pm q$.  Keep the ramification point $\infty$
as the fifth mark and omit the other ramification point.
The square-root choice $q$ and the chosen lift $+1$ give
the exact labels in (B.51).  This constructs the two
maps of spaces and therefore their path-groupoid functors.

Let $\ell_0,\ell_1$ be meridians of $0,1$ in the
$\lambda$-line, based at $q^2$.  Covering monodromy is
$\ell_0\mapsto1$ and $\ell_1\mapsto0$ in $\mathbb Z/2$.
The two-vertex lifted graph of the standard two-loop
spine has the three free generators
\[
 u_0=\ell_0^2,\qquad u_+=\ell_1,\qquad
 u_-=\ell_0\ell_1\ell_0^{-1}.
\]
Take the stems to $q=1$ along the positive real interval.
The stem to $q=-1$ first turns $q$ positively through a
half-circle around $0$ and then follows the negative
real interval.  The loop around $q=0$ rotates the inner
two points once, giving $\sigma_2^2$.  A loop around
$q=1$ simultaneously turns the two disjoint adjacent
pairs at $-1$ and $1$, giving
$\sigma_1^2\sigma_3^2$.  The chosen half-circle stem
interchanges the inner two points by $\sigma_2$; this
gives the last word in (B.52).  These observations can
also be checked directly from the four trajectories
$-1,-q,q,1$; the disjoint small meridian discs make the
two simultaneous twists commute.

The middle braid has linking number $1$ on strands
$1,2$ and the others have linking number zero.  Removing
that one full scalar braid and applying the exact combing
algorithm yields (B.53).  The program
\texttt{b\_cover\_map.py} carries out these three word
calculations and verifies the braid identities.

For the core arrows the lifted parameter path is
$q=ir$, $1/\sqrt3\leq r\leq1$, whose projection is
the real interval $\lambda=-r^2$.
For the one-node families it is exactly (B.46)--(B.47),
and for the quarter-turn it is
$q=\epsilon e^{\pi it/2}$, projecting to
$\lambda=\epsilon^2e^{\pi it}$.  Formula (B.49)
supplies the remaining composites.  This gives the
claimed expression for every arrow in the realized
subdiagram by actual four-point covering paths,
with no dimension-only identification.
\end{proof}


\subsection{Two-node plumbing and the triple-triangle families}

\begin{result}{B.24}{Proved}{A complete specified realization of six triple-triangle families}
The realization of B.22 extends by the three object families
$\mathrm{III}_2,\mathrm{III}_1,\mathrm{III}'_2$ and the six
arrow families
\[
 \tau_{\pi,\rho},\quad\gamma_{\pi,\rho},\quad
 \tau_r^1,\quad\gamma_r^1,\quad
 \gamma_r^2,\quad\gamma_r^{2'}.
\]
The explicit choices below give $422$ objects and $1020$ arrows.
There are $599$ cycles after a spanning tree and $600$ defining
relators.  The resulting comparison to $\Pi_1M_{0,5}$ is an
equivalence.  Every one of the $660$ added path images and
$420$ added relators is computed in
\texttt{b\_second\_tier.json}.  All arrows of this realization
have explicit expressions through the core paths of B.21 and
the stated twists and one-node erasures.

The triple-triangle twist is a positive quarter-turn.  The
supplied transcription's $1/3$ is not the reading used here:
the enlarged original scan 31 (folio 30), independently viewed
by two readers, supports $1/4$ and matches the upper quarter-turn
glyph.  The alternative transcription is retained in the reading
register.  The endpoint and path conventions below are part of
the realization; no additional unrecorded direct-path convention
is inferred from an endpoint count.
\end{result}
\begin{proof}
Put $\epsilon=1/10$ and label the following five-point
configurations by $0,1,2,3,4$ in the displayed order:
\[
\begin{aligned}
 R&=(\infty,-1-\epsilon,1+\epsilon,-1+\epsilon,1-\epsilon),\\
 M&=(\infty,-1-i\epsilon,1+\epsilon,-1+i\epsilon,1-\epsilon),\\
 L&=(\infty,-1-i\epsilon,1+i\epsilon,-1+i\epsilon,1-i\epsilon).
\end{aligned}                                                     \tag{B.54}
\]
The two node clusters are $A=\{1,3\}$ and $B=\{2,4\}$;
the remaining mark is $0$.  Near their centres $-1,1$ use
\[
 z_A=z+1,\quad z_B=z-1,\qquad
 z_Av_A=q_A,\quad z_Bv_B=q_B.
\]
The two three-point tails have coordinates with marks
$-1,1,\infty$.  Thus (B.54) are respectively the smoothing
parameters $(\epsilon,\epsilon)$,
$(i\epsilon,\epsilon)$ and $(i\epsilon,i\epsilon)$.
Discs of radius $1/4$ about $-1,1$ are disjoint and contain
the corresponding pairs.  Choose the two real rays on each
untwisted seam; after a quarter-turn retain their union with
the rotated two rays.  This gives the $2+2$, $4+2$ and
$4+4$ seam markings of the three diagrams.

The holomorphic stabilizers of $R,M,L$ are, respectively,
\[
 H=\langle(12)(34)\rangle,\qquad 1,\qquad H.
                                                               \tag{B.55}
\]
For $R,L$ the nonidentity element is $z\mapsto-z$.
To exclude any additional automorphism, a permutation of
the five labels determines at most one fractional linear map.
Normalize its first, second and fifth labels to $0,1,\infty$
and compare the remaining two cross-ratios.  The exact rational
calculation for all $120$ permutations gives exactly (B.55);
the finite test and all configurations are in the program.
Hence the object sets are
$\mathfrak S_5/H,\mathfrak S_5,\mathfrak S_5/H$, of sizes
$60,120,60$.

Here are all six representative paths.  In the first two,
replace the positive circular quarter-turn by its chord
\[
 u(t)=\epsilon(1-t+it),\qquad 0\leq t\leq1.
\]
The radial interpolation to the circular path avoids zero
inside the relevant disc and fixes all other marks.  The
two twist paths are
\[
\begin{aligned}
 T_A(t)&=(\infty,-1-u(t),1+\epsilon,-1+u(t),1-\epsilon),\\
 T_B(t)&=(\infty,-1-i\epsilon,1+u(t),-1+i\epsilon,1-u(t)).
\end{aligned}                                                     \tag{B.56}
\]
Their endpoints are $R\to M$ and $M\to L$.

Define the two target configurations
\[
\begin{aligned}
 R_*&=(\infty,-(1-\epsilon)^{-1},0,
                 -(1+\epsilon)^{-1},-1/2),\\
 M_*&=(\infty,-(1-i\epsilon)^{-1},0,
                 -(1+i\epsilon)^{-1},-1/2).
\end{aligned}                                                     \tag{B.57}
\]
The path $G_R$ is the straight coordinate interpolation
from $R$ to $R_*$, and $G_M$ is that from $M$ to $M_*$,
keeping the first coordinate $\infty$.  In $G_R$ the real
ordering $1<3<4<2$ is preserved.  In $G_M$ the conjugate
pair $\{1,3\}$ stays off the real line and the two other
finite marks remain real and ordered.  Thus both avoid
collisions.  More precisely the retained $A$ pair has
half-separation at most $10/99$; its centre is separated
from the other finite marks by more than $12/25$.  A moving
disc of radius $1/4$ therefore retains that pair throughout.
The complement has four marked boundary or puncture positions.
Consequently these paths are actual one-node erasures with
a fixed three-point tail, including its displayed framing.
At the endpoints use the change of coordinate $z\mapsto
-1/(z+1)$ to compare the retained disc with the one-node
model of B.22.  Its restriction transports the cut and all
its pins; a smaller disc and the intervening annulus give
the same isotopy class.

For the last two paths put $s=1-t$ and
$v=\epsilon+(1-\epsilon)t$.  Set
\[
\begin{aligned}
 G_L(t)&=(\infty,-s-i\epsilon,s+iv,-s+i\epsilon,s-iv),\\
 G_{LL}(t)&=(\infty,-1-iv,1+iv,-1+iv,1-iv).
\end{aligned}                                                     \tag{B.58}
\]
For $t<1$ in $G_L$ the two pair centres have different
real parts; at $t=1$ the imaginary magnitudes are
$\epsilon$ and $1$, so there is still no collision.
The pair $A$ remains in a disc of radius $1/4$ around
$-s$: indeed the distance from this centre to either
other point is
\[
 \sqrt{4(1-t)^2+(\epsilon+(1-\epsilon)t)^2}>1/2.
\]
Thus $G_L$ also carries a specified retained cut.  When
comparing its endpoint with the two-pin right object, use
the two pins selected by the transported normal ray;
the extra pins are explicitly erased.  This is the combined
circle and marking erasure convention for this realization.
The path $G_{LL}$ erases both cuts and reaches a square.
It fixes the real parts $-1,1$, so again no collision occurs.

To state every endpoint without an implicit relabelling,
permutations act on the labels on the left.  Write
\[
\begin{aligned}
 p_R&=(2,1,3,0,4),&
 p_L&=(0,1,3,4,2),&
 p_Q&=(0,4,3,1,2),
\end{aligned}                                                     \tag{B.59}
\]
where tuples list the images of $0,1,2,3,4$.
Then (B.57) are $p_R V(\epsilon)$ and
$p_R V(i\epsilon)$ after $z\mapsto-1/(z+1)$.
At the end of $G_L$, multiplication by $-i$ identifies
the configuration with $p_L V(\epsilon)$.  At the end
of $G_{LL}$, multiplication by $(1-i)/2$ identifies it
with $p_QK(1)$.  This proves the following endpoint table:
\[
\begin{array}{c|c|c|c|r}
 \text{family}&\text{parameter}&\text{source}&\text{target}&
 \text{arrows}\\ \hline
 \tau_{\pi,\rho}&g\in\mathfrak S_5&gH\ (R)&g\ (M)&120\\
 \gamma_{\pi,\rho}&g\in\mathfrak S_5&gH\ (R)&gp_RH\ (\mathrm{II}_2)&120\\
 \tau_r^1&g\in\mathfrak S_5&g\ (M)&gH\ (L)&120\\
 \gamma_r^1&g\in\mathfrak S_5&g\ (M)&gp_RH\ (\mathrm{II}'_2)&120\\
 \gamma_r^2&g\in\mathfrak S_5&gH\ (L)&gp_LH\ (\mathrm{II}_2)&120\\
 \gamma_r^{2'}&gH\in\mathfrak S_5/H&gH\ (L)&gp_QH_Q\ (\mathrm I_4)&60.
\end{array}                                                     \tag{B.60}
\]
In the last row $p_Q^{-1}Hp_Q\subset H_Q$, and the
path is equivariant for $H$, so its coset parameter is
well defined.  The other five rows have trivial arrow
stabilizer.  The table proves all endpoint valencies
directly and records the three permutations that would
otherwise be hidden in a comparison of diagrams.

Keep B.22's indices, then number the new objects $182$--$241$,
$242$--$361$ and $362$--$421$ in lexicographic coset or
permutation order.  Append the arrows in the order (B.60),
also lexicographically.  Normalize each of their five-point
paths by (B.41) and use the same base configuration and
least-integer detour transports as B.21.  The computation
certifies $724$ rational subintervals by exact Sturm tests.
Every pair difference stays in a strict half-plane on each
subinterval; straight interpolation therefore supplies an
actual collision-free homotopy to the polygon used for
the braid calculation.  Crossing signs use the same exact
projection of slope $1/37$.  All $660$ resulting pure-braid
words are combed and checked by the faithful Artin action.

Extend the old spanning tree greedily by these arrows.
It now has $421$ edges.  It is already connected before
the last three rows of (B.60); its new edges are twists
and $G_R$ erasures.  The graph has
$1020-422+1=599$ free cycles.  For each of its $420$ new
non-tree cycles impose its computed word in the eleven
core cycles, using B.21's inverse dictionary.  Together
with the old $180$ relators this gives $600$ relators.
All $420$ expanded equations are retained and independently
checked in the output.  Eliminating their new cycle
generators reduces the presentation exactly to B.22.
The actual path functor, its inverse on the eleven core
cycles, and the tree transports therefore prove the
claimed equivalence.

Finally, the same equations express every added arrow
as a tree transport, a word in the core cycles, and an
inverse tree transport.  The new tree transports are
the twists and the displayed one-node erasures.  The
core paths have the exact four-point covering construction
of B.23.  This gives the stated expressions through
these specified relative-dimension-one operations,
including the simultaneous two-cut erasure $G_{LL}$.
The remaining tetrahedron--triangle families, and any
further comparison with unstated choices in the manuscript,
are not supplied by this argument.
\end{proof}


\subsection{A full geometric realization and its complete relator file}

\begin{result}{B.25}{Proved}{All 542 objects and 1620 arrows in a specified realization}
There is an explicitly specified geometric realization of the
nine object families and fourteen named arrow families of
cote 148, scans 31--33, with
\[
 542\text{ objects},\qquad1620\text{ arrows},\qquad
 1079\text{ free cycles},\qquad1080\text{ defining relators}.
                                                               \tag{B.61}
\]
The corresponding path-groupoid functor to $\Pi_1M_{0,5}$
is an equivalence after imposing those relators.  All arrow
images, tree transports and relators are computed in
\texttt{b\_third\_tier.json}, using the exact conventions below
and those of B.22--B.24.  Every arrow has a computed expression
in the sixty core paths and a tree whose edges are twists or
the specified one-node erasures.  The core paths have B.23's
exact four-point covering description.

This theorem specifies the direct $f_{\vec Q,a}$ paths,
the retained-pin erasures, and the three tetrahedral endpoint
choices.  It does not identify a further manuscript homotopy
convention that is not determined by the supplied diagram.
The literal assignment question in B.15 retains precisely
that comparison, together with the general colimit of B.8.
\end{result}
\begin{proof}
Write $\zeta=(-1+i\sqrt3)/2$ and $\xi=(1+i\sqrt3)/2$.
The tetrahedron--triangle model is
\[
 W(q)=(1,\zeta,\zeta^2,-q^{-1},q^{-1}),\qquad
 W=W(\epsilon),\quad\epsilon=1/10.                \tag{B.62}
\]
At $q=0$ its stable limit has a tetrahedral component
$(1,\zeta,\zeta^2,\infty)$ and a three-point tail.
In the coordinate $w=1/z$ near that node, the two
tail marks are $-q,q$.  Thus $wv=q$ is again the normal
plumbing relation.  Exact comparison of the two normalized
cross-ratios for all $120$ permutations shows that $W$
has no nonidentity holomorphic relabelling stabilizer.
This supplies the remaining $120$ objects, indexed by
$g\in\mathfrak S_5$.

The positive sixth-turn is $q(t)=\epsilon e^{\pi it/3}$.
It is homotopic, within the punctured small disc, to
\[
 q(t)=\epsilon(1-t+\xi t).                         \tag{B.63}
\]
Throughout this homotopy the three tetrahedral marks stay
on the unit circle and the two tail marks have modulus
at least $10$, so no collision occurs.  Multiplication
of $W(\epsilon\xi)$ by $\zeta^{-1}$ identifies its endpoint
with $p_\tau W$, where
\[
 p_\tau=(1,2,0,4,3)=(012)(34).                    \tag{B.64}
\]
This is exactly the six-step permutation $\omega\times(-1)$
of the source's seam pins $J\times\varepsilon$.
Indeed the cyclic permutation of $J=\{0,1,2\}$ has
order three and the interchange of $\varepsilon=\{3,4\}$
has order two.  Its six translates concatenate to one
positive turn in the normal coordinate $q$, hence to
the separating meridian, with no suppressed full twist.
The endpoint map for this family is
\[
 \tau_r^0:g\longrightarrow gp_\tau
 \quad(g\in\mathfrak S_5).                        \tag{B.65}
\]

For $\gamma_r^0$ use the path
\[
 E(t)=(1,\zeta,\zeta^2,-(1-t)/\epsilon,
                         1/(\epsilon(1-t))),
 \qquad0\leq t\leq1,                             \tag{B.66}
\]
with the limiting value at $t=1$.  Its two moving marks
stay respectively on $(-\infty,0]$ and $[10,\infty]$,
away from each other and from the tetrahedral triple.
Its endpoint is $(1,\zeta,\zeta^2,0,\infty)$, the
bipyramid already identified with B.20's $K(1/\sqrt3)$.
The source has the conjugate pair $\{1,2\}$ and three
real marks $\{0,3,4\}$ throughout.  After fixing the
three real marks, this component is an upper or lower
half-plane.  Thus (B.66) also records its unique path
class in that real component.  Its endpoint map is
\[
 \gamma_r^0:g\longrightarrow gH_B,
                                                               \tag{B.67}
\]
so it has six incoming arrows at each of the twenty
bipyramids.

It remains to specify the three choices in the
$360$-member family $\gamma_r^{1'}$.  Use the semi-left
model $M$ of (B.54).  Normalize its real labels $0,2,4$
to $0,1,\infty$.  The point with label $1$ becomes
\[
 \lambda_*={-19+i\over181}\in\mathbb H,
\]
and label $3$ becomes its conjugate.  Therefore this
entire real component is the explicit map
\[
 C:\mathbb H\longrightarrow M_{0,5},\qquad
 C(\lambda)=(0,\lambda,1,\overline\lambda,\infty).
                                                               \tag{B.68}
\]
Its inverse takes the labelled cross-ratio determined by
$0,2,4$; thus it is a homeomorphism onto the component,
not merely a dimension count.

Put
\[
 p_0=(0,1,3,4,2),\qquad c=(0\,2\,4),\qquad
 p_j=c^jp_0\quad(j=0,1,2).                         \tag{B.69}
\]
These are the stipulated pin choices.  The tetrahedral
triple at $p_jW$ is the conjugate pair $\{1,3\}$ together
with $c^j(0)$; the other two real marks form its triangular
tail.  Cyclically permuting the three real labels preserves
the half-plane.  The three endpoint coordinates are exactly
\[
\begin{aligned}
 \lambda_0&=-\frac{30}{91}+\frac{110\sqrt3}{819}i,\\
 \lambda_1&=\frac{27}{37}+\frac{30\sqrt3}{407}i,\\
 \lambda_2&=\frac{121}{40}+\frac{33\sqrt3}{40}i.
\end{aligned}                                                     \tag{B.70}
\]
All have positive imaginary part.  Define the representative
arrow of choice $j$ by the straight segment
\[
 C((1-t)\lambda_*+t\lambda_j).                    \tag{B.71}
\]
The upper half-plane is convex, so this is collision free
and gives its unique relative path class there.  Relabelling
by $g$ supplies all arrows, with exact endpoint formula
\[
 \gamma_{r,j}^{1'}:g\ (\mathrm{III}_1)\longrightarrow
                      gp_j\ (\mathrm{II}_1),
 \qquad g\in\mathfrak S_5,\quad j=0,1,2.            \tag{B.72}
\]
There are three distinct targets at each source, and for
each target and each $j$ there is exactly one source.
This proves both valencies three, all $360$ arrows, and
the three free $\mathfrak S_5$-orbits retained in B.12.
The alternative half-plane endpoint pin choices are
not silently equated with (B.69); that selection is the
explicit convention of the present realization.

Choice $j=0$ is supported on a genuine four-point piece.
In (B.68), its three marks $0,\lambda,\overline\lambda$
stay in $|z|<1/2$, because
\[
 |\lambda_*|^2=2/181<1/4,\qquad
 |\lambda_0|^2=400/2457<1/4,
\]
and the disc is convex.  The other two marks are $1$
and $\infty$.  The original radius-$1/4$ circle at the $B$ cluster of $M$
becomes the circle of centre $-8/21$ and radius $20/21$.
Interpolate its centre to $0$ and its radius to $3/4$;
throughout, $0,\lambda_*,\bar\lambda_*$ stay inside
and $1,\infty$ stay outside.  This is an explicit isotopy
of the retained cut.  Keep the circle $|z|=3/4$: its exterior
is a fixed three-point tail, and its interior is the
varying four-point piece with its boundary.  Erasing
the other internal cut gives precisely this one-node
path.  On the retained seam uniformly refine the
two pins to six, using the orientation and the two
existing pins.  This is the uniform overmarking of
cote 148, scan 57, of multiplier three.  It supplies
the target's six pins explicitly; it is not an
identification of two different cardinalities of pin sets.
Thus the $120$ arrows of choice zero are the required
specified operations from a piece of modular dimension one.

Append the $120$ objects in lexicographic permutation
order to B.24.  Append the $120$ sixth-turns, $120$
paths (B.66), then the three $120$-arrow families
(B.71), in the order $j=0,1,2$.  This produces the
$542$ objects and $1620$ arrows of (B.61), retaining
all previous indices.

The path computation uses the same normalization,
basepoint, detour transports and projection as B.21.
For the new paths, choose the following homogeneous representatives:
for (B.63), use $(1:1),(\zeta:1),(\zeta^2:1),(-1:q),(1:q)$;
for (B.66), use $(1:1),(\zeta:1),(\zeta^2:1),
(-(1-t):\epsilon),(1:\epsilon(1-t))$; for (B.71), use
$(0:1),(\lambda:1),(1:1),(\bar\lambda:1),(1:0)$.
Their pair determinants $D_{ij}$ are polynomials of degree
at most two over $\mathbb Q(\sqrt3)(i)$.  On each rational
subinterval, every Bernstein coefficient $z$ of $D_{ij}$
satisfies
\[
 \operatorname{Re}(z\overline{D_{ij}(m)})>0,\qquad
 |\operatorname{Im}(z\overline{D_{ij}(m)})|
 <(2-\sqrt3)\operatorname{Re}(z\overline{D_{ij}(m)}),
\]
where $m$ is the midpoint.  These exact inequalities put
all determinants within $\pi/12$ of their midpoint rays.
The normalized difference of any two finite coordinates is
\[
 \frac{D_{ij}D_{ac}D_{bc}}{D_{ic}D_{jc}D_{ba}},
\]
up to a constant sign, where $a,b,c$ are the three
normalizing labels.  Its argument therefore varies by
less than $6\pi/12=\pi/2$ from its midpoint ray.  This
is the strict half-plane condition for the collision-free
homotopy to its endpoint chord.  A single determinant
certificate consequently proves the homotopy for all
$120$ relabellings.  The five representative families
need respectively $4,8,3,12,16$ subintervals, totaling
$43$.  The code retains their exact Bernstein coefficients
and all $600$ normalized braid images.  Exact crossing
arithmetic and the faithful Artin action verify every
pure-braid word.  An earlier slower run, retained as
\texttt{b\_third\_tier\_bernstein\_wip}, was stopped
before completion; no result uses that incomplete run.

For the spanning tree retain B.24's $421$ edges and
add exactly the $120$ arrows (B.72) with $j=0$.
Their targets run bijectively over the new objects,
so this is a tree with $541$ edges.  Its added edges
are precisely the one-node paths just described.
There are $1620-542+1=1079$ non-tree cycles.  For each
of the $480$ new cycles, the program records its exact
word in B.21's eleven core cycles and imposes that
one relation.  Together with the previous $600$
relations this gives $1080$ relators.  All expanded
words are retained and checked by the Artin action.
Eliminating these $480$ new generators gives exactly
B.24.  This proves completeness of the relator set,
the equivalence to $\Pi_1M_{0,5}$, and the explicit
inverse functor by core words and tree transport.

The same formula is the promised expression for each
arrow through the sixty core paths and the $541$
tree edges: its based cycle is replaced by its recorded
core word, then the source and target tree transports
are restored.  The tree edges are the specified
twists and one-node erasures, and B.23 gives the exact
four-point covering functor for the core.  Thus this
construction includes a computed relative-dimension-one
dictionary for every arrow of the specified full model.
It leaves no unspecified word in that dictionary;
the still-open issue is the explicitly retained comparison
of conventions with the manuscript.
\end{proof}

\begin{result}{B.26}{Proved}{The computed defining-relator counts are minimal}
For the connected graphs used in B.21, B.22, B.24 and
B.25, a groupoid presentation with isotropy
$\Gamma_{0,5}$ requires at least, respectively,
\[
 12,\qquad180,\qquad600,\qquad1080
\]
defining relators.  The presentations constructed there
attain these bounds.
\end{result}
\begin{proof}
The explicit split extension of B.19 is
\[
 1\longrightarrow F(x,y,z)\longrightarrow\Gamma_{0,5}
 \longrightarrow F(a,b)\longrightarrow1.
\]
Each of its six displayed conjugation relations acts
as the identity on the abelianization of the fibre.
The homology spectral sequence has only the columns
and rows numbered $0,1$, since both groups are free.
There is therefore no possible nonzero differential
affecting total degree two, and it gives
\[
 H_1(\Gamma_{0,5};\mathbb Z)=\mathbb Z^5,\qquad
 H_2(\Gamma_{0,5};\mathbb Z)
   =H_1(F_2;\mathbb Z^3)=\mathbb Z^6.              \tag{B.73}
\]
The $H_1$ equality also follows immediately from the
displayed presentation and its split projection.

Collapse a spanning tree of a connected generator graph
and write $b$ for its number of free cycles.  A set of
$r$ defining relators gives a two-dimensional presentation
complex $Y$ with one vertex, $b$ edges and $r$ two-cells.
Its first rational homology has dimension five.  The
cellular chain complex consequently gives
\[
 \dim_{\mathbb Q}H_2(Y;\mathbb Q)=r-b+5.
\]
The classifying map $Y\to B\Gamma_{0,5}$ is surjective
on $H_2$: attach cells of dimensions at least three
to kill the higher homotopy groups of $Y$; these may
quotient $H_2$ but cannot create new degree-two classes.
Equation (B.73) therefore gives $r-b+5\geq6$, or
$r\geq b+1$.  The four cycle counts are $11,179,599,1079$.
Their constructed relator counts are exactly one greater,
which proves the assertion.
\end{proof}


\subsection{An actual finite colimit of cut-surface groupoids}

\begin{result}{B.27}{Proved}{A strict finite colimit with the five elementary types}
Let $S$ be an oriented stable surface with individually labelled
boundary circles $I$, let $H=\Gamma(S,\partial S)$, and put
$A=\mathbb Z^I$, with its positive boundary-twist generators.
There is a strict finite diagram of finite-object groupoids,
with a free action of $\mu=A/2A$ on their object sets, whose
ordinary colimit is equivalent to the fundamental groupoid of
the full tangent-ray bundle of the corresponding moduli stack,
restricted to the specified finite rays.  The equivalence
preserves the boundary half-turns and full turns.

The pieces are obtained by cutting and sewing from the five
types
\[
 (0,3),\quad(0,4),\quad(0,5),\quad(1,1),\quad(1,2). \tag{B.74}
\]
For genus zero the index can be the finite poset of reduced
labelled trees $G$ with
\[
 \sum_v(\operatorname{val}(v)-3)\leq2,
\]
ordered by edge contraction.  The local five-point and elliptic
pieces here are the explicitly presented groupoids of B.19,
B.25 and B.9--B.11.  For a pinned graph with several vertices,
finite products and the flag-to-edge induction of B.1 give the
same result with the full external lattice
$\mathbb Z^{E(G)\amalg I(G)}$.

This proves a geometric finite-colimit recipe with all five
pieces (B.74).  It does not remove the elliptic pieces in favour
of only three- and four-point sphere pieces, nor identify an
unstated original finite ray convention.  Those are the stronger
comparisons retained in B.8 and B.15.
\end{result}
\begin{proof}
We give the index, all functors, the strictification and the
comparison inverse.  In particular, the colimit here
means an ordinary colimit of a specified strict diagram.

Put $d=3g-3+|I|$ and let $X=|P_2|$, with $P_2$ the
multicurves having at least $d-2$ components.  When $d\leq2$
the empty multicurve is allowed.  The simple connectedness
theorem proved and audited in D.5 applies to this precise
$P_2$.  For $d\leq2$ the asserted construction is one of
(B.74), already explicitly computed.  For $d\geq3$, every
vertex of $P_2$ is a nonempty multicurve.  Cutting a multicurve
$\sigma$ gives
\[
 \sum_Q(3g(Q)-3+n(Q))=d-|\sigma|\leq2.             \tag{B.75}
\]
Each term is nonnegative.  The stable types of nonnegative
complexity at most two are exactly (B.74).

Use the barycentric action category: its cells are flags of
multicurves in $P_2$, and its arrows are face inclusions.
Take one representative of every $H$-orbit of cells and every
incidence orbit.  The resulting index $\mathcal I$ is finite:
the dual stable graph has bounded numbers of vertices and
edges, and a flag is a chain of at most three edge subsets.
If a face has several incidences they remain separate arrows.
There is no quotient of those arrows to connected components.
Cell stabilizers are the stabilizers of the full flags, so
they fix the barycentric cells pointwise.

For a representative cell $i$, write $K_i\subset H$ for
its actual stabilizer.  It contains $A$ centrally.  Cutting
all curves of its finest multicurve describes $K_i$ using
the products of the component boundary-fixed groups, the
identifications of paired boundary twists, and the finite
automorphisms preserving the flag.  The paired twist map
is exactly the map of B.1, not a newly postulated quotient.
All component groups have the explicit presentations (B.74).
For finite graph permutations retain their actual lifts and
multiplication cocycles.  One can construct their finite
presentation by choosing one lift of each graph permutation
and adding its conjugation and multiplication relations in
the cut groups.  These relations are computed by restriction
to the labelled components and collars; the component word
problems are the braid and free-group word problems already
specified in B.9, B.18 and B.19.  This procedure makes no
splitting assertion for a finite graph-automorphism extension.
The nontrivial lift and flag calculations of B.11 and D.27
are explicit instances of this construction.

Choose actual incidence transporters $g_\alpha\in H$.
They send the chosen lifted face to its representative.
For an arrow $\alpha:i\to j$ the associated local functor
on vertex groups is
\[
 \phi_\alpha:K_i\longrightarrow K_j,\qquad
 k\longmapsto g_\alpha k g_\alpha^{-1}.            \tag{B.76}
\]
When an incidence stabilizer is required, $i$ denotes
that full flag, so (B.76) is defined on precisely that
stabilizer.  The transporter choices for a composition
give the natural isomorphism
\[
 c_{\beta,\alpha}
   =g_{\beta\alpha}(g_\beta g_\alpha)^{-1}:
 \phi_\beta\phi_\alpha\Longrightarrow\phi_{\beta\alpha}.
                                                               \tag{B.77}
\]
The element is in the target stabilizer.  Substitution
shows its naturality and its associativity identity:
both ways of comparing three composed conjugations
cancel to $g_{\gamma\beta\alpha}
(g_\gamma g_\beta g_\alpha)^{-1}$.
For identity arrows choose transporter $1$.
Thus this is an actual unital pseudofunctor, including
its specified natural transformations.

Here is its framed version.  Write $E=\tfrac12A$ and
identify $\mu=E/A$ by the indicated basis.  Define
\[
 K_i^{1/2}=(K_i\times E)/
       \{(a,-a):a\in A\},\qquad
 \chi_i:K_i^{1/2}\longrightarrow E/A.             \tag{B.78}
\]
The original $K_i$ injects and is the kernel of $\chi_i$.
Let $F(i)$ be the translation groupoid on the finite
set $\mu$: an element $h\in K_i^{1/2}$ is an arrow
$\epsilon\to\epsilon+\chi_i(h)$.  Equations (B.76)--(B.77)
extend by the identity on $E$ and hence give equivariant
functors and natural transformations on these groupoids.
The positive element $\frac12e_\ell$ gives an arrow
between opposite leg rays, and its square is exactly
$e_\ell=z_\ell$.  Translation on $\mu$ is free.

This formal description is the concrete tangent-ray
groupoid: choose a base framing and rotate each boundary
coordinate through zero or a positive half-turn.  A
path between two such framings is an element of (B.78),
and composition is concatenation.  B.1 identifies the
same construction under sewing.  If local finite models
retain additional node-ray objects, choose the positive
half-turn paths from one node-ray choice to the others
and use
\[
 f\longmapsto \ell_{t(f)}^{-1}f\,\ell_{s(f)}.       \tag{B.79}
\]
This is an equivalence to the stated stabilizer model.
For one node its first half-turn maps to $1$ and its
second to the full node twist.  Thus (B.79) retains
the meridian and does not quotient its lattice.
It is equivariant for all external leg rays.

We now replace the pseudofunctor by a strict one without
losing finiteness or freeness.  Define $\widetilde F(i)$
to have objects
\[
 (j,\alpha:j\to i,x),\qquad x\in\operatorname{Ob}F(j),
\]
and morphisms
\[
 \operatorname{Hom}_{\widetilde F(i)}
 ((j,\alpha,x),(k,\beta,y))
 =\operatorname{Hom}_{F(i)}(F(\alpha)x,F(\beta)y).
                                                               \tag{B.80}
\]
Composition is that of $F(i)$.  For $u:i\to l$, the
object map appends $u$, and the arrow map is
\[
 f\longmapsto c_{u,\beta}\,F(u)(f)\,c_{u,\alpha}^{-1}.
                                                               \tag{B.81}
\]
The associativity identity checked after (B.77) proves
that two successive maps (B.81) equal the map for their
composite, on the nose.  Thus $\widetilde F$ is a strict
finite diagram.  Evaluation into $F(i)$ is fully faithful
and essentially surjective, with inverse on objects
$x\mapsto(i,1_i,x)$ and the identity comparison arrows
from (B.80).

There is also a direct description of its ordinary colimit.
Its objects are the pairs $(i,x)$; the object
$(j,\alpha,x)$ is the image of $(j,1_j,x)$.
For each incidence $\alpha:j\to i$, the identity morphism
of $F(\alpha)x$ in (B.80) supplies an isomorphism from
$(j,x)$ to $(i,F(\alpha)x)$.  The remaining relations
are precisely the local groupoid relations, naturality
of these isomorphisms, and (B.77).  Conversely these
generators define a cone on (B.80)--(B.81), by the
displayed identity comparison arrows.  The two universal
maps fix every local morphism and every incidence
isomorphism, and are inverse.  This explicitly identifies
the ordinary colimit with the groupoid presentation of
the specified pseudofunctor.

Freeness is now checked for this particular colimit.
Every object map (B.81) preserves the tags $j,x$.
Consequently the object identifications never identify
two different elements of the sign torsor attached to
one $j$.  The colimit has the free sign action on its
objects $(j,\epsilon)$.  This is the additional argument
missing from the unrestricted colimit inference refuted
in B.5.

To verify the comparison and its inverse, let $H^{1/2}$
be (B.78) with $K_i$ replaced by $H$.  It acts on $X$
through $H/A$, so the half-boundary turns act trivially
on $X$.  Form the action-path groupoid whose objects are
$(v,\epsilon)$, with $v$ a vertex of the barycentric
subdivision of $X$, and whose arrows are
\[
 (h,[p]):(v,\epsilon)\longrightarrow(w,\epsilon+\chi(h)),
 \quad p\text{ a path from }hv\text{ to }w.
                                                               \tag{B.82}
\]
Use endpoint-preserving homotopy classes of paths.
Composition is $(k,[q])(h,[p])=(kh,[q*(kp)])$.
The local groupoids and incidence isomorphisms above
map to (B.82) by their actual transporter and cell paths.
A path in the cell complex reads conversely as a word
in those generators.  Backtracking cancels a generator
with its inverse; crossing a two-simplex gives exactly
(B.77).  This proves that the two word maps are inverse
on the full subgroupoid of the chosen orbit representatives.
It is the groupoid form of the action-complex proof,
with its generator and relation maps retained.

The functor from (B.82) to the translation groupoid of
$H^{1/2}$ sends $(h,[p])$ to $h$.  It is fully faithful:
for every $h$ the required path exists because $X$
is connected, and any two choices are homotopic because
$X$ is simply connected.  An inverse sends every sign
object to $(v_0,\epsilon)$ for a chosen reference vertex
$v_0$, and sends $h$ to $(h,[p_h])$, where $p_h$ is
any edge path from $hv_0$ to $v_0$.  The preceding
word-reading procedure computes its representative in
the finite colimit.  The choice has no effect on that
morphism, and concatenation proves the multiplication
law.  Chosen paths from the other representative
vertices to $v_0$ give the natural isomorphism from
the other composite to the identity.  In local
coordinates the forward functor sends a morphism
$f$ of (B.80) to
\[
 g_\beta^{-1}\,f\,g_\alpha
\]
in $H^{1/2}$.  This is the announced actual comparison,
including its inverse and all conjugating transports.

In genus zero, an individually leaf-labelled reduced
tree has no nonidentity automorphism: each internal
edge is determined by its two leaf sets, and the
vertices are then determined by incidence.  A graph
contraction is determined by the subset of edges retained.
Any two flags with the same labelled trees are related
by a surface homeomorphism preserving those edge
partitions.  Hence the orbit flags are exactly the
chains in the finite tree poset stated in the theorem.
Moreover the stabilizer of a finer tree fixes every
edge separately, so it preserves every contracted face.
We may therefore use that poset itself with (B.76)--(B.77);
its barycentric action presentation is the same one
just proved.  Strictification (B.80) does not change
its index.

Finally, for a pinned graph with several vertices,
apply the construction at each vertex, take the finite
product diagram, and apply the central induction and
ray-torsor tensor map of B.1 to each value and arrow.
Products of groupoids preserve colimits separately,
as follows from their adjunction with functor groupoids.
Central induction preserves them by its defining
extension-of-operators adjunction; the explicit
flag-to-edge map adds the two exponents at a sewn node.
Thus the resulting colimit is the actual sewn groupoid,
and its comparison is the sewn product of the already
specified functors.  Reapplying the sign-object
construction with the external lattice
$\mathbb Z^{E(G)\amalg I(G)}$ gives the same freeness
proof and the required node and leg half-turns.
Finite graph symmetries can be retained as an action
groupoid on these pinned objects; this introduces
arrows and never takes their coarse object quotient.
\end{proof}


\subsection{An explicit sphere sewing dictionary}

\begin{result}{B.28}{Proved}{Finite pants and chamber presentations compared on every generator}
After capping the leg framings, the finite chamber groupoids
of B.18 for $M_{0,4}$ and $M_{0,5}$ have the following
finite pants presentations and exact inverse dictionaries:
\[
\begin{array}{c|r|r|r|r}
 &\text{objects}&F\text{ arrows}&B\text{ arrows}&\text{relators}\\ \hline
 M_{0,4}&12&6&12&7\\
 M_{0,5}&120&120&144&169 .
\end{array}                                                     \tag{B.83}
\]
Here each $B$ is a positive half-turn of a three-boundary
piece, and each $F$ changes one cut in a four-boundary
piece.  The $24$ associativity relators in the second
row are genuine five-point pentagons.  Both comparison
functors and their inverse natural isomorphism are
given below.  The program \texttt{b\_parenthesized.py}
computes every arrow and relation word.

Consequently every based braid word computed in B.21--B.25
has an explicit word in sphere $F$ and $B$ moves by the
displayed substitution.  This proves the sewing dictionary
for those capped based words.  It retains the separate
questions of the original finite ray transports and the
elliptic framed $S$-move coherence.
\end{result}
\begin{proof}
Put $m=n-1$.  The objects are pairs $(p,T)$, where
$p\in\mathfrak S_m$ is an ordering of the finite labelled
points on the real line and $T$ is a full binary
parenthesization of their ordered positions.  Use the
outer boundary as the $n$th leg.  The two shapes for
$m=3$, and the five for $m=4$, give the object counts.

These are concrete pants markings.  Around each
consecutive block in $T$ take a small rounded rectangle
containing its real points; choose the collar width
strictly decreasing with the depth of the block.
Parent rectangles contain child rectangles, and sibling
rectangles are disjoint.  Their boundaries give the
specified pants cuts.  Moving the widths through these
strict inequalities gives an isotopy, so no additional
mapping class is hidden in the choice.

There is a directed $F$ arrow for every occurrence
\[
 ((A\,B)\,C)\longrightarrow(A\,(B\,C)).
\]
It forgets the first internal cut and inserts the second
inside the same sphere with four boundary components.
Its underlying punctured surface map is the identity.
This is the inverse-removal followed by removal convention
already compared in B.14.  For every cherry consisting
of consecutive individual leaves $i,i+1$, include a
positive $B$ arrow interchanging those two points in
their pair disc.  It fixes the outside of that disc.
The endpoint ordering is $ps_i$, and its underlying
braid is exactly $\sigma_i$.

Define the forward functor $\Phi$ to the chamber model
by
\[
 \Phi(p,T)=p,\qquad
 \Phi(F)=1_p,\qquad
 \Phi(B_{p,T,i})=(p,\sigma_i).                    \tag{B.84}
\]
To give its inverse, choose the left-associated shape
$L_m$.  In the finite $F$ graph for each ordering,
choose the breadth-first tree rooted at $L_m$, with
the lexicographic shape and arrow order used in the
program.  Let $q_{p,T}$ be its unique path from
$(p,L_m)$ to $(p,T)$.  For each $i$, let $T_i$ be the
first shape in the program's order having the cherry
$i,i+1$.  The inverse on the chamber generator is
\[
 \Psi(p,\sigma_i)=
 q_{p,T_i}\, B_{p,T_i,i}\,q_{ps_i,T_i}^{-1}
                                                               \tag{B.85}
\]
in traversal order, and $\Psi(p)=(p,L_m)$.
The image of (B.85) under (B.84) is literally $\sigma_i$,
including its source and target.

Here is the complete finite relation list, rather than
an appeal to an unspecified presentation.  First, for
every non-tree $F$ edge $e:(p,T)\to(p,U)$ impose
\[
 q_{p,T}\,e\,q_{p,U}^{-1}=1.
                                                               \tag{B.86}
\]
Second, for every $B$ edge that is not the chosen
representative at $T_i$, impose
\[
 q_{p,T} B_{p,T,i}q_{ps_i,T}^{-1}
   =q_{p,T_i}B_{p,T_i,i}q_{ps_i,T_i}^{-1}.
                                                               \tag{B.87}
\]
Third, substitute (B.85) into every lifted Artin
relation of B.18, and into its one positive full-scalar
relation.  The latter caps the outer boundary as in
B.18.  All these words and their typed endpoints are
listed in \texttt{b\_parenthesized.json}.

Relations (B.86) say exactly that every $F$ path is
its tree path.  When $m=4$, the five-shape $F$ graph
is a pentagon, so there is one such relator for each
of the $24$ label orderings.  Its geometric equality
is B.14's cut-removal pentagon.  Relation (B.87)
compares two actual realizations of the same pair-disc
half-turn; both have the same braid $\sigma_i$.
The remaining Artin relations are the usual disjoint
support and three-strand isotopies.  The program checks
every resulting braid equality by the faithful Artin
action; the last relation maps to the full scalar braid,
whose capping quotient was proved in B.18.

For completeness, contract the $F$ trees.  Relations
(B.86) remove every remaining $F$ generator, and
(B.87) remove all duplicate $B$ generators.  What
remains is exactly B.18's chamber presentation, on
its labelled vertices, arrows and relators.  Thus
$\Psi$ is defined on the presented groupoid.
The composite $\Phi\Psi$ is the identity on every
object and generator.  The paths $q_{p,T}$ give
the natural isomorphism from $\Psi\Phi$ to the
identity: its naturality for $F$ is (B.86), and
for $B$ is (B.87).  This proves the claimed
equivalence with its actual inverse maps.

The finite enumeration gives $6$ $F$ and $12$ $B$
arrows at $m=3$.  There are no relations (B.86) or
(B.87); the six lifted braid relations and the
scalar relation give seven.  At $m=4$ there are
$120$ $F$ and $144$ $B$ arrows, $24$ relations
(B.86), $72$ relations (B.87), $72$ lifted Artin
relations and one scalar relation, totaling $169$.
These counts and all maps were computed and checked.

Finally, take a computed based word from B.21--B.25.
Read its successive adjacent braid crossings with
their changing label order and replace each positive
crossing by (B.85), each negative crossing by the
inverse word with its changed source.  Equation
(B.84) recovers the original word exactly, so this
is the promised explicit sphere sewing factorization.
The non-braid tree transports of the original finite
ray diagrams remain the stated transports of B.25;
no assertion about an unrecorded framing convention
is obtained by dropping them.
\end{proof}
